\documentclass[a4paper,11pt,oneside,BCOR=1cm,DIV=12,toc=chapterentrywithdots]{scrbook}

\usepackage[T1]{fontenc}
\usepackage[utf8]{inputenc}
\usepackage{graphicx}
\usepackage{booktabs}
\usepackage{tabularx}
\usepackage{array}
\usepackage{amsmath}
\usepackage{enumitem}
\usepackage{xcolor}
\usepackage[expansion=false]{microtype}
\usepackage{listings}
\usepackage[nohyperlinks]{acronym}
\usepackage{scrlayer-scrpage}
\usepackage{hyperref}
\usepackage{xstring}
\setlength{\emergencystretch}{2em}

\definecolor{tucgreen}{RGB}{0,120,90}
\definecolor{codegray}{RGB}{245,247,249}
\hypersetup{
  colorlinks=true,
  linkcolor=tucgreen,
  citecolor=tucgreen,
  urlcolor=tucgreen,
  bookmarksnumbered=true,
  bookmarksopen=true,
  pdfauthor={David Keci},
  pdftitle={Extending Jupyter Notebook Reproducibility Analysis Beyond GitHub}
}

\clearpairofpagestyles
\ihead{Cross-Platform Notebook Reproducibility}
\ohead{Draft}
\cfoot{\pagemark}
\setkomafont{pageheadfoot}{\small}
\setkomafont{chapter}{\color{tucgreen}}
\setkomafont{section}{\color{tucgreen}}

\lstset{
  basicstyle=\ttfamily\small,
  backgroundcolor=\color{codegray},
  frame=single,
  rulecolor=\color{black!15},
  breaklines=true,
  columns=fullflexible,
  keepspaces=true,
  showstringspaces=false,
  aboveskip=1em,
  belowskip=1em
}

\newcommand{\systemname}{Notebook FAIR Cross-Platform Pipeline}
\newcommand{\draftwarning}{\textbf{Draft status:} Chapters 1--7 document the completed cross-platform pipeline, notebook classifier, and knowledge-graph extension. Evaluation, discussion, and conclusion follow the planned thesis structure.}

\begin{document}
\hypersetup{pageanchor=false}

\begin{titlepage}
  \begin{center}
    \raisebox{-1ex}{\includegraphics[scale=1.5]{../template/TU_Chemnitz_positiv_gruen.pdf}}\\[1.2cm]
    {\Large\bfseries Extending Jupyter Notebook Reproducibility Analysis Beyond GitHub\\[0.35cm]}
    {\large A Cross-Platform Pipeline for Codeberg and Zenodo\\[1.3cm]}
    {\Large\bfseries Master Thesis - Draft}\\[0.8cm]
    Submitted in Fulfilment of the Requirements for the Academic Degree\\
    Master of Science\\[0.5cm]
    Department of Computer Science\\
    Chair of Distributed and Self-organizing Systems
  \end{center}
  \vfill
  \begin{tabular}{@{}ll@{}}
    Submitted by: & David Keci\\
    Programme: & M.Sc. Web Engineering\\
    Supervisor: & Prof. Dr. Sheeba Samuel\\
    University: & Chemnitz University of Technology\\
    Draft date: & 18 August 2026
  \end{tabular}
\end{titlepage}

\hypersetup{pageanchor=true}
\frontmatter

\addchap{Abstract}
Jupyter notebooks combine executable code, narrative text, and generated results, but their apparent self-containment does not guarantee computational reproducibility. Existing large-scale reproducibility pipelines and the downstream FAIR Jupyter knowledge graph primarily operate on GitHub-hosted notebooks. This focus excludes resources distributed through alternative Git forges and research repositories whose access mechanisms and metadata models differ from GitHub. This thesis develops a cross-platform extension that processes GitHub, Codeberg, and Zenodo resources through a common reproducibility route while preserving the distinctions between Git repositories and archived research records.

The implementation introduces platform detection, platform-specific validation and acquisition, normalized metadata extraction, mixed-platform batch processing, and a copy-on-first-run database architecture. GitHub and Codeberg resources are validated and acquired with Git, whereas Zenodo records are validated through the Records API and downloaded as archives. Eight descriptive metadata fields are normalized and stored in a dedicated table. After acquisition, the existing platform-independent stages collect notebooks, infer requirements, create Python environments, execute notebooks, compare outputs, and record provenance and error information.

The database is subsequently exported into seven CSV datasets and transformed into RDF using YARRRML-derived RML mappings and Morph-KGC. An ontology extension represents platforms, Git repositories, archived records, notebooks, repository runs, notebook executions, reproducibility results, and notebook-classification activities. The confirmed reproducibility snapshot generated 222,048 unique triples and passed structural and SPARQL competency checks. A hybrid classifier compares transparent weighted rules with a structured AI-model response, stores both decisions, and requests human clarification when the methods disagree or have low confidence.

\vspace{0.6cm}
\noindent\textbf{Keywords:} computational reproducibility, Jupyter notebooks, metadata normalization, knowledge graphs, GitHub, Codeberg, Zenodo

\vspace{0.6cm}
\noindent\draftwarning

\addchap*{Table of Contents}
{\color{tucgreen}\hrule height 0.7pt}
\vspace{0.45cm}

\newcommand{\tocpagerange}[1]{\StrBehind{#1}{. }[\tocpagerangevalue]\tocpagerangevalue}
\newcommand{\tocchapterline}[3]{\noindent\textcolor{tucgreen}{\textbf{#1\quad #2}}\dotfill\textbf{\tocpagerange{#3}}\par\vspace{0.28em}}
\newcommand{\tocsectionline}[3]{\noindent\hspace{1.5em}\textcolor{black!82}{#1\quad #2}\dotfill\textcolor{black!65}{\tocpagerange{#3}}\par}
\newcommand{\tocsubsectionline}[3]{\noindent\hspace{3em}\textcolor{black!68}{#1\quad #2}\dotfill\textcolor{black!55}{\tocpagerange{#3}}\par}
\newcommand{\tocfrontline}[2]{\noindent\textcolor{tucgreen}{\textbf{#1}}\dotfill\textbf{#2}\par\vspace{0.28em}}

\begingroup
\small
\setlength{\parskip}{0.08em}

\tocfrontline{Abstract}{1 page}

\tocchapterline{1.}{Introduction}{pp. 1--4}
\tocsectionline{1.1}{Motivation}{p. 1}
\tocsectionline{1.2}{Problem Statement}{p. 2}
\tocsectionline{1.3}{Research Goal and Research Questions}{pp. 2--3}
\tocsectionline{1.4}{Contributions}{p. 3}
\tocsectionline{1.5}{Thesis Structure}{p. 4}

\tocchapterline{2.}{Background and Related Work}{pp. 5--13}
\tocsectionline{2.1}{Computational Reproducibility}{p. 5}
\tocsectionline{2.2}{Jupyter Notebooks and Reproducibility}{p. 6}
\tocsectionline{2.3}{Existing Notebook Reproducibility Analysis}{p. 7}
\tocsectionline{2.4}{GitHub, Codeberg, and Zenodo}{p. 8}
\tocsectionline{2.5}{Metadata Normalization and Provenance}{pp. 8--9}
\tocsectionline{2.6}{Environment Reconstruction and Notebook Execution}{p. 10}
\tocsectionline{2.7}{Notebook Classification Approaches}{p. 11}
\tocsectionline{2.8}{FAIR Data and Knowledge Graphs}{p. 12}
\tocsectionline{2.9}{Research Gap}{p. 13}

\tocchapterline{3.}{Requirements and System Analysis}{pp. 14--19}
\tocsectionline{3.1}{Baseline System and Scope}{p. 14}
\tocsectionline{3.2}{Functional Requirements}{p. 15}
\tocsectionline{3.3}{Non-Functional Requirements}{p. 16}
\tocsectionline{3.4}{Input and Output Contracts}{p. 17}
\tocsectionline{3.5}{Data and Classification Requirements}{p. 18}
\tocsectionline{3.6}{Evaluation Requirements}{p. 19}

\tocchapterline{4.}{System Design}{pp. 20--27}
\tocsectionline{4.1}{Overall Architecture}{pp. 20--21}
\tocsectionline{4.2}{Pipeline Control Flow}{p. 22}
\tocsectionline{4.3}{Connector Architecture}{p. 23}
\tocsectionline{4.4}{Metadata Model and Normalization Design}{p. 24}
\tocsectionline{4.5}{Persistence and Database Design}{p. 25}
\tocsectionline{4.6}{Notebook Classification Architecture}{p. 26}
\tocsectionline{4.7}{Knowledge Graph and Provenance Design}{p. 27}

\clearpage
\tocchapterline{5.}{Cross-Platform Pipeline Implementation}{pp. 28--42}
\tocsectionline{5.1}{Project Structure and Pipeline Entry Point}{p. 28}
\tocsectionline{5.2}{Run Creation, Single and Batch Processing}{p. 29}
\tocsectionline{5.3}{Platform Detection and Validation}{pp. 30--31}
\tocsectionline{5.4}{GitHub Baseline Integration}{p. 32}
\tocsectionline{5.5}{Codeberg Connector Implementation}{pp. 33--34}
\tocsectionline{5.6}{Zenodo Connector and Content Acquisition}{pp. 35--36}
\tocsectionline{5.7}{Cross-Platform Metadata Normalization}{p. 37}
\tocsectionline{5.8}{Notebook Registration and Discovery}{p. 38}
\tocsectionline{5.9}{Dependency and Environment Reconstruction}{pp. 39--40}
\tocsectionline{5.10}{Notebook Execution and Output Comparison}{p. 41}
\tocsectionline{5.11}{Error Handling, Testing, and Container Verification}{p. 42}

\tocchapterline{6.}{Notebook Classification}{pp. 43--51}
\tocsectionline{6.1}{Classification Objective and Taxonomy}{p. 43}
\tocsectionline{6.2}{Feature Extraction and Input Representation}{p. 44}
\tocsectionline{6.3}{Rule-Based Notebook Classification}{pp. 45--47}
\tocsubsectionline{6.3.1}{Classification Rules}{p. 45}
\tocsubsectionline{6.3.2}{Rule Priorities and Conflict Resolution}{p. 46}
\tocsubsectionline{6.3.3}{Rule-Based Output}{p. 47}
\tocsectionline{6.4}{AI-Based Notebook Classification}{pp. 48--49}
\tocsubsectionline{6.4.1}{Model Selection}{p. 48}
\tocsubsectionline{6.4.2}{Prompt Design}{pp. 48--49}
\tocsubsectionline{6.4.3}{Input Context Construction}{p. 49}
\tocsubsectionline{6.4.4}{Structured AI Output}{p. 49}
\tocsectionline{6.5}{Hybrid Classification Strategy}{p. 50}
\tocsectionline{6.6}{Classification Integration into the Pipeline}{pp. 50--51}
\tocsectionline{6.7}{Failure Handling and Classification Examples}{p. 51}

\tocchapterline{7.}{Knowledge Graph Extension}{pp. 52--61}
\tocsectionline{7.1}{Relationship to FAIR Jupyter Knowledge Graph}{p. 52}
\tocsectionline{7.2}{Relational Data Export}{p. 53}
\tocsectionline{7.3}{Ontology Extension}{pp. 54--55}
\tocsectionline{7.4}{YARRRML and RML Mapping Design}{p. 56}
\tocsectionline{7.5}{Knowledge Graph Materialisation}{pp. 57--58}
\tocsectionline{7.6}{Representation of Notebook Classifications}{p. 59}
\tocsectionline{7.7}{SPARQL Competency Questions}{p. 60}
\tocsectionline{7.8}{Resulting Knowledge Graph}{p. 61}

\clearpage
\tocchapterline{8.}{Evaluation}{pp. 62--73}
\tocsectionline{8.1}{Evaluation Strategy and Dataset}{p. 62}
\tocsectionline{8.2}{Connector and Metadata Correctness}{p. 63}
\tocsectionline{8.3}{Cross-Platform Reproducibility Experiment}{pp. 64--65}
\tocsectionline{8.4}{Cross-Platform Notebook Characteristics}{p. 66}
\tocsectionline{8.5}{Rule-Based Classification Evaluation}{p. 67}
\tocsectionline{8.6}{AI-Based Classification Evaluation}{p. 68}
\tocsectionline{8.7}{Hybrid Classification Evaluation}{p. 69}
\tocsectionline{8.8}{Knowledge Graph Validation}{pp. 70--71}
\tocsectionline{8.9}{End-to-End and Integration Testing}{p. 72}
\tocsectionline{8.10}{Error Analysis}{p. 73}

\tocchapterline{9.}{Discussion}{pp. 74--77}
\tocsectionline{9.1}{Answering the Research Questions}{p. 74}
\tocsectionline{9.2}{Design Decisions and Trade-offs}{p. 75}
\tocsectionline{9.3}{Threats to Validity}{p. 76}
\tocsectionline{9.4}{Ethical and Operational Considerations}{p. 77}

\tocchapterline{10.}{Conclusion and Future Work}{pp. 78--80}
\tocsectionline{10.1}{Summary of Contributions}{p. 78}
\tocsectionline{10.2}{Conclusion}{p. 79}
\tocsectionline{10.3}{Future Work}{p. 80}

\tocfrontline{Bibliography}{}
\tocfrontline{Appendices}{}
\endgroup

\mainmatter
\chapter{Introduction}
\label{chap:introduction}

\section{Motivation}
\label{sec:motivation}

Computational results are reproducible when an independent party can take the same data, code, and computational steps and obtain consistent results~\cite{nasem2019}. Reproducibility has become a practical concern in computational research because an increasing share of published findings depends on software that must still run correctly years after publication. Jupyter notebooks~\cite{kluyver2016} are a central artefact in this setting. They interleave executable code, narrative explanation, and stored outputs in a single document~\cite{nbformat}, which makes them attractive for communicating an analysis. The same property also creates a false impression of self-containment: a notebook shows its code and its results, but it does not necessarily record the library versions, the Python interpreter, the input data, or the cell order that produced those results.

Empirical studies confirm that this gap is large rather than marginal. Pimentel et al.~\cite{pimentel2019} analysed more than 1.1 million notebooks collected from GitHub and re-executed a substantial subset under controlled conditions; only a minority finished without errors, and fewer still reproduced their stored outputs. Samuel and Mietchen~\cite{samuel2024computational} repeated this kind of analysis for notebooks associated with biomedical publications and reported a similar pattern, with missing dependencies and unavailable data dominating the failure causes. Related work reinforces the picture from other angles: notebooks frequently contain duplicated and re-ordered code~\cite{koenzen2020duplication}, large-scale corpus studies find pervasive quality problems~\cite{psallidas2022datascience}, and automated environment restoration remains an unsolved problem in the general case~\cite{wang2020restoring}. Reproducibility of published notebooks is therefore an open, measurable research problem rather than a solved engineering detail.

\subsection*{Where Notebooks Are Published}

Researchers do not publish computational results in one place. Code is typically shared through a Git forge, while data, software snapshots, and supplementary materials are increasingly deposited in an archival research repository that issues a persistent identifier. The resulting landscape is heterogeneous, and its components differ substantially in size, governance, access mechanism, and metadata model.

GitHub is by a large margin the dominant Git forge. Its 2024 Octoverse report states more than 518 million total projects and more than 1.5 million repositories that use Jupyter notebooks, a figure that grew by roughly 92\% in a single year~\cite{octoverse2024}. Codeberg is a much smaller, non-profit alternative operated by Codeberg e.V. in Berlin and built on the Forgejo forge software~\cite{codebergdocs,forgejoapi}; public figures report on the order of 300{,}000 repositories and 200{,}000 registered accounts~\cite{codebergstats}, roughly three orders of magnitude smaller than GitHub. GitLab occupies an intermediate position and exposes yet another API model~\cite{gitlabdocs}. On the archival side, Zenodo, operated by CERN, held approximately 7.3 million records at the time of writing~\cite{zenodostats}, each carrying a DOI, descriptive metadata, and downloadable files~\cite{zenodoapi}; Figshare~\cite{figshare} serves a comparable role for other communities. Table~\ref{tab:platform-landscape} summarises this landscape.

\begin{table}[ht]
\centering
\caption{Publication venues for computational notebooks. Figures are approximate and reflect the sources cited at the time of writing.}
\label{tab:platform-landscape}
\small
\begin{tabularx}{\textwidth}{@{}p{1.9cm}p{2.5cm}p{2.9cm}X@{}}
\toprule
Platform & Kind & Approximate scale & Primary access mechanism\tabularnewline
\midrule
GitHub & Git forge & $>$518M projects; $>$1.5M notebook repositories~\cite{octoverse2024} & Git clone plus GitHub REST API~\cite{githubapi}\tabularnewline
GitLab & Git forge & Large; self-hosted instances widespread & Git clone plus GitLab REST API~\cite{gitlabdocs}\tabularnewline
Codeberg & Git forge (non-profit) & $\approx$300K repositories; $\approx$200K accounts~\cite{codebergstats} & Git clone plus Forgejo REST API~\cite{forgejoapi}\tabularnewline
Zenodo & Archival repository & $\approx$7.3M records~\cite{zenodostats} & Records API; files downloaded as archives~\cite{zenodoapi}\tabularnewline
Figshare & Archival repository & Millions of research outputs~\cite{figshare} & REST API; files downloaded as archives\tabularnewline
\bottomrule
\end{tabularx}
\end{table}

\subsection*{Why Non-GitHub Venues Matter for Reproducibility Research}

The evidence base described above is drawn almost entirely from GitHub. Pimentel et al.~\cite{pimentel2019}, Samuel and Mietchen~\cite{samuel2024computational}, and the FAIR Jupyter knowledge graph~\cite{samuel2024fairjupyter} that publishes the latter study's results all sample GitHub repositories. This is a defensible choice given GitHub's scale, but it leaves an unexamined question: what is currently known about notebook reproducibility is what is true of one platform, one community of practice, and one metadata model.

There are three concrete reasons why this limitation matters. First, publishing practice is not uniform across venues. An archival Zenodo deposit is created at a specific moment, keeps its deposited files under a persistent identifier, and is usually made to accompany a publication; a Git repository is a living working copy whose state changes after publication. These are different kinds of artefact, and there is no a priori reason for their reproducibility properties to coincide. Second, metadata quality is platform-dependent by construction. Zenodo requires a title, creators, and a publication date before a DOI is minted, whereas a Git forge requires almost nothing beyond a repository name. If metadata completeness correlates with reproducibility, that correlation cannot be detected from a single-platform sample. Third, the incentive structures differ: a non-profit forge such as Codeberg attracts projects with different motivations than a commercial one, and a repository deposited to satisfy a journal's data-availability policy is produced under different constraints than a personal working repository.

Answering whether these differences matter requires infrastructure that does not currently exist. It requires a workflow that can acquire, describe, execute, and compare notebooks from heterogeneous venues on identical terms, while keeping the origin of every result visible. Building that workflow is not merely a matter of accepting additional hostnames.

\section{Problem Statement}
\label{sec:problem-statement}

The reproducibility pipeline of Samuel and Mietchen~\cite{samuel2024computational} provides a working route from a repository to structured execution evidence: it acquires the repository, infers dependencies, reconstructs a Python environment, executes the notebooks, compares the new outputs against the stored outputs, and records the outcome in a relational database. The pipeline~\cite{samuel2024computational} assumes throughout that every input is a GitHub repository. That assumption is embedded in its URL handling, its validation step, its acquisition step, and its database lookup. Extending the analysis beyond GitHub raises four distinct problems.

\paragraph{Heterogeneous acquisition.} A Git forge and an archival repository expose fundamentally different objects. Codeberg is superficially similar to GitHub --- both serve Git repositories --- but it runs Forgejo, so its API base, JSON field names, and web notebook links (\texttt{/src/} rather than \texttt{/blob/}) differ~\cite{forgejoapi}. Zenodo is different in kind: a record is not a Git remote at all, so it must be validated through the Records API, and its files must be downloaded and extracted before any notebook exists locally~\cite{zenodoapi}. Substituting a different clone URL does not address either case.

\paragraph{Metadata heterogeneity.} Each platform describes a resource with its own vocabulary. GitHub reports a single owner login, Zenodo reports a list of creators; a licence may be an SPDX identifier, a nested object, or absent entirely; a DOI exists on Zenodo and normally does not on a Git forge. Without a shared representation, equivalent fields cannot be compared across platforms. With a careless shared representation, the distinction between a genuinely missing value and a value that the pipeline invented is lost --- and it is precisely the missing values that carry the interesting signal about metadata quality.

\paragraph{Provenance across a heterogeneous corpus.} Provenance is the record of the entities, activities, and agents involved in producing a result~\cite{prov2013}. Once results originate from several platforms, provenance stops being optional bookkeeping: an execution outcome is interpretable only in combination with the platform it came from, the resource it belongs to, the processing run that produced it, and the notebook it describes. The relational schema and the downstream knowledge graph must therefore carry this chain explicitly.

\paragraph{Confounding by notebook purpose.} Execution success is not a property of the platform alone. A tutorial notebook, a data-cleaning notebook, and a machine-learning experiment have different dependency footprints, different runtimes, and different failure modes. If one platform happens to host proportionally more tutorials, a naive comparison of success rates would report a platform effect where a composition effect exists. Controlling for this requires that notebook purpose be characterised systematically, and the existing pipeline does not characterise it at all.

\section{Research Goal and Research Questions}
\label{sec:research-questions}

The goal of this thesis is to investigate whether and how computational-notebook reproducibility can be studied across heterogeneous research platforms, and to determine what such a cross-platform view reveals that a single-platform view cannot. This goal is decomposed into four research questions.

\begin{description}[leftmargin=1.3cm,labelwidth=1.05cm,style=nextline]
  \item[\textbf{RQ1}] To what extent can a platform-independent reproducibility workflow support the analysis of computational notebooks across heterogeneous research platforms?
  \item[\textbf{RQ2}] How does metadata heterogeneity across research platforms affect the integration and representation of computational notebooks and their reproducibility-related information?
  \item[\textbf{RQ3}] To what extent do rule-based and AI-based approaches agree when characterizing computational notebooks into reproducibility-relevant categories, and to what extent do the two approaches complement each other?
  \item[\textbf{RQ4}] To what extent do computational notebooks differ across platforms with respect to metadata completeness, execution outcomes, failure causes, and reproducibility-relevant categories?
\end{description}

RQ1 asks whether platform-independent processing generalises: the interesting question is not whether connectors can be written, but at which point in the workflow platform differences genuinely stop mattering and where they persist. RQ2 concerns the representational consequences of heterogeneity --- what is lost, what is preserved, and what becomes measurable when descriptions from incompatible metadata models are integrated. RQ3 treats notebook characterisation as a measurement problem and asks whether two independent automated approaches converge on the same description of a notebook, and whether their respective failures are independent enough to be worth combining. RQ4 uses the resulting instrument to ask the empirical question that motivated the work, without presupposing that any platform is more reproducible than another. Section~\ref{sec:eval-strategy} maps each research question to the evidence that addresses it, and Section~\ref{sec:rq-answers} states the answers reached.

\section{Contributions}
\label{sec:contributions}

This thesis contributes a platform-independent instrument for studying computational-notebook reproducibility, together with an empirical comparison of reproducibility evidence across a commercial Git forge, a non-profit Git forge, and an archival research repository. In detail:

\begin{itemize}[leftmargin=1.4em]
  \item \textbf{A platform-independent reproducibility workflow.} A connector abstraction isolates platform-specific validation, description, and acquisition behind a common contract, after which notebook discovery, dependency inference, environment reconstruction, execution, and output comparison proceed identically regardless of origin. The abstraction is demonstrated for a commercial forge, a Forgejo-based forge, and an archival repository, and is designed so that a further venue requires a connector rather than a second pipeline.
  \item \textbf{A normalized cross-platform metadata model.} Eight descriptive fields are reconciled across three incompatible API vocabularies with an explicit non-fabrication rule: a value the source does not provide is stored as absent, never inferred. This turns metadata incompleteness from an artefact of processing into a measurable property of the platform, and makes the comparison in RQ4 possible.
  \item \textbf{An evidence-preserving relational design.} Stable entities are separated from processing activities, so that repeated attempts accumulate as history rather than overwriting it, the inherited baseline database is never modified, and every execution outcome remains traceable to its platform, resource, run, and notebook.
  \item \textbf{A hybrid, auditable notebook characterisation method.} A transparent weighted-rule classifier and a locally hosted AI model independently label the same bounded notebook evidence. Both decisions, their confidences, their agreement state, and the resulting review requirement are retained, so that the reliability of the characterisation can itself be evaluated rather than assumed.
  \item \textbf{A semantic publication layer extending FAIR Jupyter.} The FAIR Jupyter ontology~\cite{samuel2024fairjupyter} is extended with hosting platforms, archived research records, pipeline runs, normalized metadata, and classification activities, and the graph is materialised reproducibly from the database through versioned YARRRML/RML mappings~\cite{heyvaert2018yarrrml,rmlspec} and Morph-KGC~\cite{arenas2024morph}. The result is queryable with SPARQL~\cite{sparql2013} and rebuildable on demand.
  \item \textbf{An empirical cross-platform evaluation.} A stratified corpus spanning all three platforms is processed end to end, and metadata completeness, run outcomes, execution outcomes, failure causes, and cell-level output agreement are reported per platform, with counts alongside every percentage.
\end{itemize}

Together these contributions provide a system that supports reproducibility analysis over heterogeneous research infrastructure, and an evidence base that quantifies how far reproducibility findings established on one platform transfer to others. Chapter~\ref{chap:design} develops the design in detail, and Chapter~\ref{chap:evaluation} reports what the instrument measured.

\section{Thesis Structure}
\label{sec:thesis-structure}

The remainder of the thesis is organized as follows.

\begin{description}[leftmargin=2.25cm,labelwidth=1.95cm,style=multiline]
  \item[\textbf{Chapter 2}] describes the terminology used throughout the thesis --- computational reproducibility, notebook structure, metadata normalization, provenance, environment reconstruction, and knowledge graphs --- and discusses the state of the art in notebook reproducibility analysis, concluding with the research gap that this work addresses.
  \item[\textbf{Chapter 3}] analyses the baseline system and derives the functional, non-functional, data, and evaluation requirements that a cross-platform reproducibility workflow must satisfy, each stated so that it can be verified.
  \item[\textbf{Chapter 4}] develops the system design: the layered architecture, the control flow through a single resource, the connector abstraction, the normalized metadata model, the relational schema, the classification architecture, and the provenance and knowledge-graph model.
  \item[\textbf{Chapter 5}] describes the implementation, following the execution flow through the cross-platform pipeline, the hybrid notebook classifier, and the knowledge-graph extension, and explaining the decisions taken at each stage.
  \item[\textbf{Chapter 6}] evaluates the system: the technical validation of the connectors, metadata layer, and knowledge graph, followed by the cross-platform experiment and the assessment of the classifier, and closing with the answers to the four research questions.
  \item[\textbf{Chapter 7}] discusses the limitations of the study and the work that remains, distinguishing threats to the validity of the present results from extensions that would broaden their scope.
  \item[\textbf{Chapter 8}] concludes by summarising what was built, what was measured, and what the results imply for reproducibility research beyond a single platform.
\end{description}

\chapter{Background and Related Work}
\label{chap:background}

This chapter introduces the ideas needed to understand the system without describing its implementation. It first explains computational reproducibility and the special properties of Jupyter notebooks. It then reviews earlier notebook studies, the three source platforms, metadata and provenance, environment reconstruction, classification, and semantic publication. The final section identifies the gap addressed by this thesis.

\section{Computational Reproducibility}

Computational results are reproducible when another person can use the same input data, code, computational steps, and analysis conditions and obtain consistent results. This thesis follows the terminology of the National Academies, which separates \emph{reproducibility} from \emph{replicability}: replication answers the same research question with newly collected data, whereas reproduction repeats the original computation \cite{nasem2019}. The distinction matters because this work evaluates digital research objects. It does not repeat the scientific study or judge whether its conclusion is correct.

Reproduction is not a single yes-or-no event. A rerun can stop while the software environment is being created, fail during execution, finish with changed outputs, or finish with outputs that match the stored notebook. Each outcome provides different evidence. A missing package points to incomplete dependency information. A network error may show that an external service is no longer available. A changed numerical result may arise from a newer library, random behaviour, or updated data. Therefore, a useful pipeline records the stage, error, duration, and comparison result instead of reducing every attempt to ``failed.''

Three groups of information influence computational reproducibility. The first is the \emph{computational object}: source code, notebook cells, input files, and stored outputs. The second is the \emph{environment}: programming-language version, libraries, operating-system assumptions, and external tools. The third is the \emph{provenance}: where the resource came from, when it was retrieved, what process acted on it, and which result that process produced. Access to code alone covers only part of this evidence.

Large notebook studies confirm the practical importance of these distinctions. Pimentel et al.~\cite{pimentel2019} examined quality and reproducibility across a large notebook corpus and connected successful execution with practices such as ordered cells and dependency information. Samuel and Mietchen~\cite{samuel2024computational} analysed notebooks linked to biomedical publications and recorded environment construction, execution exceptions, and output comparison separately. Their work shows that a pipeline must preserve intermediate outcomes if the results are to support later analysis.

Two operational definitions are used throughout this thesis. \emph{Successful execution} means that every selected code cell completes in the reconstructed environment without raising an uncaught exception. \emph{Output reproducibility} is stricter and is assessed per cell: the outputs of the freshly executed notebook are compared against the outputs stored in the published notebook, and each code cell is recorded as producing an identical result, a different result, or a result recognised as nondeterministic (for example a timestamp, a memory address, or a value derived from an unseeded random source). This cell-level comparison is the measurement introduced by Samuel and Mietchen~\cite{samuel2024computational}, and it is retained unchanged in this work; the per-notebook reproducibility score reported in Chapter~\ref{chap:evaluation} is the proportion of comparable code cells whose outputs matched.

Keeping the granularity at the cell level matters. A notebook in which one of forty cells diverges is a materially different outcome from a notebook in which none of the forty matched, and a binary reproduced/not-reproduced verdict cannot distinguish the two. These definitions make the evaluation measurable while acknowledging that identical output is not proof of scientific validity: they answer the narrower question of whether the published computation can still be rerun, and how much of its recorded result survives that rerun.

\clearpage
\section{Jupyter Notebooks and Reproducibility}

A Jupyter notebook is both a document and an executable program. It can place source code beside natural-language explanations, equations, tables, figures, and interactive results. This combination makes notebooks useful for exploring data and explaining an analysis to readers \cite{kluyver2016}. It also creates a challenge: the visible document does not necessarily contain everything needed to repeat the computation.

The standard \texttt{.ipynb} file is a JSON document containing an ordered list of cells and notebook-level metadata \cite{nbformat}. Code cells may store source code, an execution count, and several output types. Markdown cells store explanatory text, and raw cells can preserve content for later conversion. Kernel and language information may appear in metadata, but most metadata fields are optional. A valid notebook can therefore omit the exact Python or package versions used by its author.

Table~\ref{tab:notebook-anatomy} lists the notebook information and its value for reproduction.

\begin{table}[ht]
\centering
\caption{Notebook information and its value for reproduction.}
\label{tab:notebook-anatomy}
\begin{tabularx}{\textwidth}{@{}lXX@{}}
\toprule
Element & What it contains & Why it matters\tabularnewline
\midrule
Code cell & Program statements and imports & Reveals operations and undeclared packages\tabularnewline
Markdown cell & Explanations, headings, and instructions & May describe purpose, data, and execution steps\tabularnewline
Stored output & Text, images, tables, or errors & Provides a reference for comparison\tabularnewline
Execution count & The recorded order of code execution & Can reveal out-of-order interactive work\tabularnewline
Notebook metadata & Kernel and language hints & Helps select a compatible environment\tabularnewline
Repository files & Requirements, setup files, and data & Supplies context missing from the notebook\tabularnewline
\bottomrule
\end{tabularx}
\end{table}

Interactive execution is the central difficulty. A user may run cell 10, return to cell 3, redefine a variable, and then save the notebook. The displayed outputs can depend on hidden kernel state that is not reproduced by a clean top-to-bottom run. Rule et al.~\cite{rule2018} explain that notebooks are often used first for trying ideas and exploring data. During this work, users may run cells in different orders or test several versions of an analysis. The saved notebook may therefore not clearly show how the final result was produced. Kery et al.~\cite{kery2018} found that data scientists often organize and explain the notebook only after the exploratory work is complete.

Reproducible execution therefore requires more than opening the file. The notebook should be run in a new kernel, in a known order, with explicit dependencies and available data. Stored outputs should be preserved long enough to compare them with new outputs. The surrounding repository is equally important because it may contain local modules, requirements files, setup scripts, or documentation. This thesis treats the notebook as the central executable object while using its repository or archived record as the context needed to interpret and rerun it.

\clearpage
\section{Existing Notebook Reproducibility Analysis}

Earlier studies examine notebook reproducibility in different ways. Some check the quality of many notebooks. Others try to execute notebooks automatically or organize the results in a knowledge graph. These studies provide the starting point for this thesis.

Pimentel et al.~\cite{pimentel2019} studied more than one million notebooks from GitHub and reran a selected group under controlled conditions. They found that reproducibility is related to features such as a clear execution order, version control, and documented dependencies. Samuel and Mietchen~\cite{samuel2024computational} later studied notebooks connected to biomedical publications. Their pipeline found the GitHub repositories, created Python environments, executed the notebooks, recorded errors, and compared the new outputs with the saved outputs. The results were stored as structured database records, not only as downloaded files.

FAIR Jupyter then converted this database information into a knowledge graph \cite{samuel2024fairjupyter}. This allows users to ask connected questions about publications, repositories, notebooks, executions, and results. However, both the source pipeline and the first graph version mainly describe GitHub resources. Other research gives practical advice to notebook authors. Rule et al.~\cite{rule2019tenrules}, for example, recommend a clear cell order, documented dependencies, and visible data-processing steps.

Table~\ref{tab:related-work} summarises the relationship between key prior work and this thesis.

\begin{table}[ht]
\centering
\caption{Relationship between key prior work and this thesis.}
\label{tab:related-work}
\small
\begin{tabularx}{\textwidth}{@{}p{2.6cm}XXX@{}}
\toprule
Work & Main focus & Useful foundation & Remaining limitation\tabularnewline
\midrule
Pimentel et al.~\cite{pimentel2019} & Large-scale notebook quality and reruns & Quality indicators and execution evidence & GitHub corpus; no cross-platform connector model\tabularnewline
Samuel and Mietchen~\cite{samuel2024computational} & Biomedical notebook reproduction & Baseline database, environment, errors, comparisons & Resource acquisition centred on GitHub\tabularnewline
FAIR Jupyter~\cite{samuel2024fairjupyter} & Semantic exploration of results & Ontology, mappings, provenance-oriented queries & GitHub repository representation\tabularnewline
This thesis & Cross-platform reproducibility & Reuses execution and semantic foundations & ---\tabularnewline
\bottomrule
\end{tabularx}
\end{table}

This thesis builds on these works~\cite{pimentel2019,samuel2024computational,samuel2024fairjupyter} instead of replacing them. The existing environment, execution, and comparison stages already record useful information and are therefore reused. The new work starts when a resource enters the pipeline. The system must now identify, check, describe, and download resources from three platforms. It must also record the platform in the database, classification results, and knowledge graph. This makes it possible to compare metadata and reproducibility across platforms, which a GitHub-only dataset cannot do.

\clearpage
\section{GitHub, Codeberg, and Zenodo}

GitHub, Codeberg, and Zenodo can all publish notebooks, but they store different kinds of resources. GitHub and Codeberg are Git hosting platforms. Their main resource is a version-controlled repository. Zenodo is a research repository. Its main resource is a published record with metadata and downloadable files. The pipeline must keep this difference because one download method cannot correctly handle all three platforms.

GitHub provides two forms of access. Git is used to clone the files, while the REST API provides information such as the name, owner, description, topics, licence, and dates \cite{githubapi}. Codeberg also supports Git, but it uses the Forgejo software and API \cite{codebergdocs,forgejoapi}. Therefore, GitHub and Codeberg can be cloned in a similar way, but their API addresses and JSON fields are not always the same~\cite{codebergapi}. Their notebook links also differ: GitHub normally uses \texttt{/blob/}, while Codeberg uses \texttt{/src/}.

Zenodo gives each published record a numeric identifier and, on publication, a DOI~\cite{zenodoapi}. Its Records API returns research metadata and a list of files that can be downloaded \cite{zenodoapi}. A record may contain notebooks, source-code archives, data, documentation, or several of these. It normally has no Git remote. Therefore, the Zenodo part of the pipeline must check the API record, select a useful file, download it, extract it, and then search the extracted folder for notebooks.

Table~\ref{tab:platform-differences} shows the platform properties that affect reproducibility processing.

\begin{table}[ht]
\centering
\caption{Platform properties that affect reproducibility processing.}
\label{tab:platform-differences}
\small
\begin{tabularx}{\textwidth}{@{}lXXX@{}}
\toprule
Aspect & GitHub & Codeberg & Zenodo\tabularnewline
\midrule
Primary object & Git repository & Git repository & Archived research record\tabularnewline
Stable local acquisition & Clone or pull & Clone or pull & Download and extract\tabularnewline
Validation signal & Git remote & Git remote & Records API\tabularnewline
Metadata source & GitHub REST API & Forgejo REST API & Zenodo Records API\tabularnewline
Identifier & owner/repository & owner/repository & numeric record ID and DOI\tabularnewline
Version context & commits and branches & commits and branches & record versions and deposited files\tabularnewline
\bottomrule
\end{tabularx}
\end{table}

These differences motivate separating platform-specific acquisition from platform-independent notebook processing. Only the early steps depend on the platform: confirming that a resource exists, reading its description, and obtaining its files. Once those files are available locally, the notebooks they contain can be discovered, executed, and compared in the same way regardless of where they came from. A study spanning several platforms therefore does not require a separate analysis workflow for each of them. The connector architecture adopted in this thesis is presented in Section~\ref{sec:connector-design}.

\section{Metadata Normalization and Provenance}

Metadata and provenance answer two different questions. Metadata answers, ``What is this resource?'' It includes information such as the title, authors, description, licence, keywords, DOI, and dates. Provenance answers, ``Where did the resource come from, and what did the pipeline do with it?'' For example, metadata can say that a Zenodo record has a certain title and two authors. Provenance can show that the record came from Zenodo, was downloaded by the pipeline, and was processed during a specific run.

GitHub, Codeberg, and Zenodo do not return this information in the same form. GitHub may provide one owner account, while Zenodo may provide a list of creators. One platform may return a licence as a short name, another may return a larger data object, and sometimes no licence is available. These differences make direct comparison difficult.

Metadata normalization is the general technique for resolving such differences. It maps the differing responses of several sources onto one shared set of descriptive elements, so that values with the same meaning can be compared regardless of where they came from. Which elements a study selects depends on the questions it asks, but established metadata standards provide a common starting point: the DataCite schema, for example, defines a small mandatory core of title, creator, publisher, publication year, and identifier for describing research objects \cite{datacite45}.

A normalization scheme has to decide what to do when a source provides no value for an element. The decision matters more than it first appears. A scheme that fills the gap, by substituting a default or by deriving a value from elsewhere, produces a uniformly complete description in which a genuinely absent value can no longer be distinguished from a supplied one. A scheme that records the absence keeps that distinction, and thereby makes metadata completeness itself measurable. For a study comparing how well different platforms describe what they host, the second choice is the only workable one.

Dates require the same care. A date reported by a source describes the resource, recording when its author created or last changed it. A date recorded by an analysis describes the analysis, recording when the resource was retrieved or processed. The two belong to different histories, and conflating them would make the record of the study indistinguishable from the record of the object being studied.

The knowledge graph records this history using PROV-O, a standard vocabulary for provenance \cite{prov2013}. It has three main ideas. An \texttt{Entity} is a resource, such as a repository, notebook, or result. An \texttt{Activity} is an action, such as a repository run or notebook execution. An \texttt{Agent} is responsible for a resource or action, such as GitHub, Codeberg, or Zenodo. For example, Zenodo is the agent that provides a record, the record is an entity used by the pipeline run, and the run is an activity that produces an execution result.

Together, normalization and provenance make results from heterogeneous sources both comparable and checkable. Normalization makes equivalent values comparable; provenance keeps visible the path by which each value was obtained, so that a surprising result can be traced back to its origin rather than merely accepted or discarded. The concrete normalization model adopted in this thesis is presented in Section~\ref{sec:metadata-model-design}, and the provenance model in Section~\ref{sec:kg-design}.

\clearpage
\section{Environment Reconstruction and Notebook Execution}

A notebook needs more than the code stored in its cells. It may depend on a specific Python version, external packages, local modules, system programs, input data, environment variables, or online services. Environment reconstruction means using the available files and notebook content to create an environment in which the notebook can run.

Files such as \texttt{requirements.txt}, \texttt{setup.py}, and Python-version files provide the clearest dependency information. When these files are missing or incomplete, imports in notebook cells provide additional clues. However, an imported module can have a different name from the package installed with \texttt{pip}, and an import normally gives no version. Automatically generated requirements are therefore only the pipeline's best attempt to rebuild the environment. They are not an exact copy of the author's original computer.

Selecting an appropriate interpreter version is equally important. Code written for one Python version can fail under another, and many packages are only compatible with a bounded version range. Version-management tools such as \texttt{pyenv} allow several interpreter releases to coexist on one machine and to be selected per project \cite{pyenv}, and installing each project's packages into a separate virtual environment prevents the dependencies of one resource from perturbing another. Wang et al.~\cite{wang2020restoring} show that automatically restoring a notebook's original execution environment is itself a hard problem, and that even careful reconstruction frequently fails to recover the author's setting exactly.

Jupyter \texttt{nbconvert} executes notebooks automatically and saves the new outputs \cite{nbconvert}. Each run starts with a new kernel and executes the cells from top to bottom. This removes variables left over from an earlier interactive session. It can also reveal hidden assumptions. For example, a later cell may need a variable that the author created by running cells in a different order. Chattopadhyay et al.~\cite{chattopadhyay2020} found that environment setup, dependency management, and turning exploratory notebooks into reliable workflows are common problems.

Execution success and output reproducibility are different results. A notebook can finish without errors but still produce different numbers, text, tables, or figures. This can happen because of random behaviour, current dates, changed online data, or new package versions. The comparison stage records these changes instead of treating every difference as a simple error. Matching outputs also do not prove that the scientific method is correct. They only show that the tested run produced the same saved result.

For this reason, a reproducibility study must treat re-execution as a staged process and observe each stage separately rather than collapsing an attempt into a single success-or-failure verdict. Dependency resolution, interpreter selection, package installation, execution, and output comparison can each fail for different reasons and carry different evidential meaning. Studies that preserve these intermediate outcomes \cite{pimentel2019,samuel2024computational} can attribute a failure to a specific cause, whereas studies that record only a final verdict cannot.

\clearpage
\section{Notebook Classification Approaches}

Reproducibility results are easier to understand when notebooks are grouped by their main purpose. A tutorial, a data-cleaning notebook, and a machine-learning experiment can need different packages and may fail for different reasons. Classification helps separate these differences from differences caused by the hosting platform.

A notebook does not normally contain one field that clearly states its purpose. Useful clues appear in the filename, headings, Markdown explanations, imports, function calls, outputs, and order of cells. Earlier studies show that notebooks often mix exploration with later explanation \cite{rule2018,kery2018}. The classifier must therefore examine both written text and code. It must also allow a notebook to have more than one important purpose.

Table~\ref{tab:classification-approaches} gives the classification strategies considered in this thesis.

\begin{table}[ht]
\centering
\caption{Classification strategies considered in this thesis.}
\label{tab:classification-approaches}
\small
\begin{tabularx}{\textwidth}{@{}p{2.7cm}XXX@{}}
\toprule
Approach & Strength & Limitation & Research relevance\tabularnewline
\midrule
Weighted rules & Transparent and repeatable; works without a model & Keywords and imports can be ambiguous & Transparent baseline\tabularnewline
AI model & Can combine code and narrative context & Output may vary and requires validation & Context-sensitive alternative\tabularnewline
Human review & Can resolve context and mixed intent & Slow and not suitable for every item & Reference for ambiguous cases and validation\tabularnewline
Hybrid decision & Exposes agreement and uncertainty & Requires a clear conflict policy & Supports analysis of agreement and uncertainty\tabularnewline
\bottomrule
\end{tabularx}
\end{table}

The categories used for such a characterisation are not arbitrary. Corpus studies of notebook content converge on a recurring set of activities. Rule et al.~\cite{rule2018} distinguish exploratory analysis from explanatory and instructional material. Kery et al.~\cite{kery2018} describe an iterative cycle in which data is first prepared, then analysed, and only later documented. Pimentel et al.~\cite{pimentel2019} separate notebooks by their dominant library usage, which cleanly distinguishes data-manipulation work from modelling work and from plotting work. Psallidas et al.~\cite{psallidas2022datascience}, analysing millions of GitHub notebooks, report the same broad activity groups --- data acquisition and preparation, exploration and analysis, visualization, and model training --- as the dominant uses of the format. Quaranta et al.~\cite{quaranta2021kgtorrent} observe an equivalent distribution in a large Kaggle corpus, and Koenzen et al.~\cite{koenzen2020duplication} additionally identify tutorial and template notebooks as a distinct population that is copied and adapted rather than authored from scratch.

The taxonomy adopted in this thesis is derived from these sources and is defined in Section~\ref{sec:classification-taxonomy}. Two observations from the literature constrain any such scheme. First, a notebook frequently serves several purposes at once, because preparation, analysis, and visualization commonly occur within one document \cite{rule2018,kery2018}; a usable scheme therefore requires an explicit convention for multi-purpose notebooks rather than forcing a single label. Second, the evidence available in a notebook is sometimes too sparse to support any confident assignment, so a scheme also requires a way to express that the evidence is insufficient instead of guessing.

These limitations motivate combining complementary automated approaches rather than relying on a single one, and reserving human review for the cases where the automated evidence is weak or contradictory. Research on human involvement in machine learning supports reviewing unclear cases instead of accepting every model answer as correct \cite{mosqueira2023hitl}. Whichever approaches are combined, their reliability cannot be established from their agreement with one another: two methods can agree and both be wrong, so a comparison against labels assigned by a person remains necessary. The classification architecture adopted in this thesis is presented in Section~\ref{sec:classification-design}.

\clearpage
\section{FAIR Data and Knowledge Graphs}

The FAIR principles state that digital research objects should be findable, accessible, interoperable, and reusable \cite{wilkinson2016fair}. They apply to data, code, workflows, and other research materials. FAIR information should be understandable not only to people but also to software. For notebook reproducibility, providing a download link is not enough. The notebook should be connected to its platform, metadata, processing history, execution, and result in a machine-readable form.

RDF stores information as simple statements called triples. Each triple has a subject, a relation, and an object \cite{rdf2014}. For example, ``repository 42 -- is hosted on -- Codeberg'' is one triple. ``run 17 -- processed -- repository 42'' is another. Because both statements use the same repository identifier, they become connected in the graph. SPARQL is a query language that follows these connections \cite{sparql2013}. It can move from a platform to its repositories, notebooks, executions, and results without manually joining several CSV files.

An ontology defines the types of things and the relations allowed in the graph. FAIR Jupyter reuses established vocabularies, among them PROV-O for processing history and DOAP for software projects, and adds terms specific to notebook reproducibility \cite{samuel2024fairjupyter}. Its published graph describes publications, GitHub repositories, notebooks, executions, and reproducibility results drawn from the biomedical study. Its vocabulary is correspondingly scoped: it models the Git repository as the unit of publication, and has no terms for a hosting platform as such, for an archived research record, or for the processing activity that produced a given result. Any study spanning more than one kind of publication venue therefore requires an extension of this model.

Graphs of this kind are normally not written by hand. Declarative mapping languages describe how records in an existing source --- a relational table or a CSV export --- become RDF statements. RML generalises the relational R2RML standard to heterogeneous sources \cite{rmlspec}, and YARRRML provides a human-readable surface syntax that compiles into RML \cite{heyvaert2018yarrrml}. A materialisation engine such as Morph-KGC \cite{arenas2024morph} then reads the mappings and emits the triples. The significance of this arrangement for reproducibility research is that the ontology, the mapping rules, and the resulting graph remain three separate, independently inspectable artefacts: a reviewer can examine the model and the transformation rules, and the graph can be regenerated whenever the underlying data changes.

A materialised graph also has to be validated. Validation normally combines structural checks --- that the ontology and mappings parse, and that source records yield the expected entities and relations --- with competency questions, that is, SPARQL queries formulating the questions the graph was built to answer. A graph that parses but cannot answer its competency questions represents a large file of triples rather than a usable model.

\clearpage
\section{Research Gap}
\label{sec:research-gap}

The work reviewed in this chapter establishes that notebook re-execution can be automated at scale \cite{pimentel2019,samuel2024computational}, that its outcomes can be recorded at cell-level granularity \cite{samuel2024computational}, and that the resulting evidence can be published semantically and queried \cite{samuel2024fairjupyter}. What it does not establish is whether any of these findings hold beyond the single platform on which they were obtained. Four gaps follow, each corresponding to one research question stated in Section~\ref{sec:research-questions}.

\paragraph{Gap 1: platform-independence has been assumed rather than demonstrated.} The execution, comparison, and error-classification stages of existing pipelines are plausibly independent of where a notebook came from, but this has never been tested, because every published corpus originates from GitHub. The boundary at which platform-specific handling genuinely ends is therefore unknown, as is whether the invariant holds for an archival record --- an object that is not a Git repository and that may package its notebooks inside an archive without any accompanying environment specification. Establishing where that boundary lies is the subject of RQ1.

\paragraph{Gap 2: the representational cost of metadata heterogeneity is unmeasured.} GitHub, Forgejo, and Zenodo describe a resource through incompatible vocabularies with different mandatory fields \cite{githubapi,forgejoapi,zenodoapi}. Integrating them requires a normalized representation, but normalization is lossy, and the direction of the loss matters: a scheme that silently fills absent values destroys exactly the signal that distinguishes a well-described archival deposit from a bare repository. No existing study reports how much descriptive information survives such an integration, nor how completeness varies by venue. This is the subject of RQ2. The FAIR Jupyter ontology \cite{samuel2024fairjupyter} provides a suitable semantic foundation but models GitHub repositories only; it has no vocabulary for a hosting platform, an archived research record, a pipeline run, or a normalized description.

\paragraph{Gap 3: notebook purpose is unmeasured, and the available automated methods have not been compared.} The corpus studies discussed in Section~2.7 agree on the broad activity types that notebooks serve, but none attaches a purpose label to the individual notebooks whose reproducibility is being measured. Two families of automated labelling are available: transparent rules over imports and text, and language models that read the notebook as a document. They rest on different evidence and can be expected to fail in different circumstances, yet how far they converge on the same description, and whether their weaknesses are independent enough to make combining them worthwhile, has not been reported. This is the subject of RQ3.

\paragraph{Gap 4: cross-platform differences are unknown and currently unmeasurable.} Because Gaps 1--3 are open, the empirical question cannot presently be asked: it is not known whether metadata completeness, execution outcomes, failure causes, or output agreement differ between a commercial forge, a non-profit forge, and an archival repository. A comparison drawn without controlling for notebook purpose would additionally risk attributing a composition effect to the platform. Closing this gap requires the instrument that Gaps 1--3 describe, and is the subject of RQ4.

These four gaps are addressed by one system rather than four separate artefacts, because they are interdependent: the empirical comparison requires the integrated metadata, which requires the platform-independent workflow, and its interpretation requires validated purpose labels. Chapter~\ref{chap:requirements} turns these gaps into verifiable requirements, Chapter~\ref{chap:design} presents the design that satisfies them, Chapter~\ref{chap:implementation} describes its realisation, and Chapter~\ref{chap:evaluation} reports the resulting evidence.

The scope of what such an instrument can settle is bounded, and the bound should be stated here rather than discovered later. A sample drawn from three platforms cannot represent the full population of notebooks on any of them. A reconstructed environment is an approximation of the author's original setting, not a recovery of it \cite{wang2020restoring}. An automated purpose label is a prediction, not a ground truth. What the system can provide is a consistent, traceable, and repeatable procedure that applies identical treatment to resources of heterogeneous origin, and that keeps the origin of every measurement visible in the result.

\chapter{Requirements and System Analysis}
\label{chap:requirements}

This chapter defines what the system must do and how each requirement can be tested.

\section{Baseline System and Scope}

The starting point is an existing pipeline for Python notebooks stored on GitHub. It reads repository and notebook information from SQLite, finds dependencies, creates a separate Python environment, executes notebooks, compares new and saved outputs, and stores the results. FAIR Jupyter exports this database information and uses it to build an RDF knowledge graph. This thesis extends those components instead of creating a separate system.

The process starts with a resource URL or a resource already listed in the database. It ends with database results, notebook categories, and an RDF graph that can be rebuilt. The extension must recognize GitHub, Codeberg, and Zenodo; download their content; convert their metadata into common fields; execute Python notebooks through the same main route; and record where each resource came from. It must also classify notebook purpose so the results can be compared by platform and category.

\begin{table}[H]
\centering
\caption{Scope of the thesis system.}
\label{tab:scope}
\footnotesize
\begin{tabularx}{\textwidth}{@{}XX@{}}
\toprule
Included & Outside the present scope\tabularnewline
\midrule
Public GitHub and Codeberg repositories & Private repositories requiring user-specific access\tabularnewline
Public Zenodo records with downloadable files & Arbitrary research repositories beyond Zenodo\tabularnewline
Python notebook environment reconstruction & Guaranteed reproduction of every operating system\tabularnewline
Execution and stored-output comparison & Validation of the scientific conclusion\tabularnewline
Normalized descriptive metadata & Automatic repair of missing source metadata\tabularnewline
Rule-based and AI-model classification & Treating automated labels as human ground truth\tabularnewline
RDF materialisation and SPARQL validation & Operation of a permanent public triple-store service\tabularnewline
\bottomrule
\end{tabularx}
\end{table}

The original database is kept read-only because it contains the starting research data. The pipeline copies it to \texttt{data/output/db/} and writes all new columns and results to that copy. Existing GitHub rows must still work after the database structure is extended.

A notebook that runs successfully is not automatically scientifically correct. Also, a limited sample cannot represent every resource on a platform. The evaluation must therefore explain how examples were selected and show the number of resources used for each result.

\clearpage
\section{Functional Requirements}

Functional requirements describe actions that the system must perform. Each requirement below can be checked with a command, database query, generated file, or test result. The requirements describe behaviour instead of individual function names, so they remain valid if the code is reorganized later.

\begin{table}[H]
\centering
\caption{Functional requirements of the extended pipeline.}
\label{tab:functional-requirements}
\small
\begin{tabularx}{\textwidth}{@{}p{0.9cm}p{9.5cm}p{2.2cm}@{}}
\toprule
ID & Requirement & Verification\tabularnewline
\midrule
FR1 & Detect GitHub, Codeberg, and Zenodo from a supported URL and reject unknown hosts. & Unit examples\tabularnewline
FR2 & Validate Git resources through Git and Zenodo resources through its record API before metadata is stored. & Connector test\tabularnewline
FR3 & Acquire Git repositories by clone or update and Zenodo records by file download and extraction. & Local files\tabularnewline
FR4 & Translate each supported API response into the same eight metadata fields without inventing missing values. & DB assertion\tabularnewline
FR5 & Support both one-resource input and platform-aware batch processing from SQLite. & End-to-end run\tabularnewline
FR6 & Register supplied notebook paths and discover \texttt{.ipynb} files automatically when paths are absent. & Notebook rows\tabularnewline
FR7 & Infer requirements, select Python, create an isolated environment, and execute Python notebooks. & Execution rows\tabularnewline
FR8 & Compare original and executed outputs and record status, metrics, errors, and elapsed time. & Result rows\tabularnewline
FR9 & Classify notebook purpose with rules and an AI model, store both attempts, and expose disagreement for review. & Classifier test\tabularnewline
FR10 & Export relational tables, materialise RDF with versioned mappings, and answer defined SPARQL questions. & KG test suite\tabularnewline
\bottomrule
\end{tabularx}
\end{table}

FR2 must happen before FR4. The pipeline checks that the source is valid before saving its metadata. Otherwise, information taken from an error response could be saved as if it were real metadata. The metadata is then fetched before the files are downloaded because the API already provides the source description and origin.

FR7 and FR8 reuse the original environment and execution steps. Each new connector only has to provide a local folder and notebook paths. It does not create a separate environment system. FR9 can run even when notebook execution fails because classification reads the notebook content without executing it. This allows failed notebooks to keep a category for later error analysis. FR10 requires the graph to be generated from the working database instead of edited by hand.

Together, these requirements describe one complete route from an input URL to results that can be queried. Downloading a file alone does not complete a connector: the resource must also be checked and described. Creating an N-Triples file alone does not complete the graph: its relationships and SPARQL test questions must also be correct.

\clearpage
\section{Non-Functional Requirements}

Non-functional requirements define the quality conditions under which the functional requirements must be fulfilled. They do not add new pipeline steps. Instead, they state how the pipeline must behave while collecting external resources, reading platform APIs, modifying database copies, executing notebooks, and publishing RDF results.

For this thesis, the most important quality goals are data safety, comparability, traceability, repeatability, robustness, security, and maintainability. These goals are necessary because a successful demonstration is not enough: the system must also protect the original data, explain how a result was produced, and remain understandable when new platforms or notebook samples are added.

\begin{table}[H]
\centering
\caption{Non-functional requirements of the extended pipeline.}
\label{tab:nonfunctional-requirements}
\footnotesize
\begin{tabularx}{\textwidth}{@{}>{\raggedright\arraybackslash}p{1.05cm}>{\raggedright\arraybackslash}p{2.45cm}>{\raggedright\arraybackslash}p{5.1cm}>{\raggedright\arraybackslash}X@{}}
\toprule
ID & Quality goal & Meaning in this thesis & Verification evidence\tabularnewline
\midrule
NFR1 & Data safety & Keep the source database unchanged and write new results only to the working copy. & Source database checksum or timestamp does not change.\tabularnewline
NFR2 & Comparability & Store platform results with shared fields, statuses, and metric definitions. & One query groups results by platform and status.\tabularnewline
NFR3 & Traceability & Link each result to platform, resource, run, execution, and classification method. & Foreign keys and RDF links connect the steps.\tabularnewline
NFR4 & Repeatability & Rebuild CSV and RDF outputs from the database, ontology, and mappings. & Generated files can be deleted and recreated.\tabularnewline
NFR5 & Robustness & Record expected failures without stopping the whole batch. & Failed resources keep status rows; later resources continue.\tabularnewline
NFR6 & Security & Treat URLs, archives, API text, tokens, and notebooks as untrusted. & Tokens are not logged; text is escaped; execution is isolated.\tabularnewline
NFR7 & Maintainability & Keep platform logic in connectors and notebook processing in shared stages. & Connector changes do not rewrite execution or RDF export.\tabularnewline
\bottomrule
\end{tabularx}
\end{table}

NFR1 protects the research baseline. The inherited GitHub database is input evidence, so the extension may read it but must not overwrite it. The first pipeline run creates a working copy in \texttt{data/output/db/}; all schema migrations, metadata rows, execution rows, and classification rows are then written to that copy. This makes experiments reversible because a new run can start again from the unchanged source database.

NFR2 and NFR3 make the results interpretable. Comparability does not mean that the platforms are treated as identical. It means that equivalent information, such as title, licence, run status, and output-comparison metrics, is stored in the same structure. At the same time, provenance keeps important differences visible: Zenodo remains an archived record, GitHub and Codeberg remain Git repositories, and every result can be followed from platform to resource, run, notebook execution, metric row, and classification activity.

NFR4 separates source evidence from generated artefacts. The working database, YARRRML/RML mappings, and ontology are the sources for semantic publication. CSV exports, N-Triples fragments, the combined graph, and SPARQL result files are generated outputs. They should therefore be rebuildable instead of manually edited. This requirement supports later verification because reviewers can inspect both the source rules and the generated graph.

NFR5 ensures that failures remain useful evidence. A batch may contain an unreachable repository, a missing Zenodo file, a package-installation failure, or a notebook execution error. These cases should be stored with a clear status and message so that error analysis can include them. They should not erase earlier results or prevent later resources in the batch from being processed.

NFR6 limits the risks created by external input. API tokens are read from environment variables and must not be written to logs or database fields. Text returned by APIs is escaped before database insertion. Downloaded archives are extracted only inside the intended resource directory. Notebook execution uses separate Python environments to avoid dependency conflicts. This isolation improves operational safety, but it is not a complete security sandbox; a public service would require stronger filesystem and network restrictions.

NFR7 keeps the design understandable. Platform connectors handle validation, metadata access, identifier normalization, and acquisition. After a local resource directory exists, notebook discovery, environment reconstruction, execution, comparison, classification, and knowledge-graph export follow shared code paths. This makes the thesis easier to evaluate because platform differences and platform-independent processing are described separately.

\section{Input and Output Contracts}

An input or output contract defines the exact information passed between parts of the system. In this pipeline, the contract is simple: the beginning of the pipeline may receive a resource in two different ways, but both ways must create the same six internal fields before notebook processing starts.

\noindent\textbf{Single-resource mode.} The input is one URL typed by a person. The URL points to a GitHub repository, a Codeberg repository, or a Zenodo record. The output is the same URL split into the six fields shown in Table~\ref{tab:io-contract}.

\noindent\textbf{Batch mode.} The input is one resource row already stored in the working database. No URL is typed by a person. The output is again the same six fields, rebuilt from the row's saved platform, identifier, and path values.

The important point is that the rest of the pipeline does not need to know which mode was used. It receives a URL, normalized identifier, platform, notebook paths, setup paths, and requirement paths. If some paths are unknown, their fields stay empty. This is common for Zenodo, because notebook paths are often only known after the record has been downloaded and extracted.

A GitHub or Codeberg notebook link needs one extra conversion. For example, a GitHub link ending in \texttt{/blob/main/analysis.ipynb} is a browser link, not a repository-relative path. The pipeline removes the display part, such as \texttt{/blob/main/} or \texttt{/src/branch/main/}, and keeps only the notebook path, for example \texttt{analysis.ipynb}. A Zenodo link is different: it identifies a record first, and the notebook paths are discovered later from the downloaded files.

\begin{table}[H]
\centering
\caption{Six internal fields produced by both input modes.}
\label{tab:io-contract}
\small
\begin{tabularx}{\textwidth}{@{}>{\raggedright\arraybackslash}p{3.2cm}>{\raggedright\arraybackslash}X>{\raggedright\arraybackslash}p{4.1cm}@{}}
\toprule
Field & Meaning & Example or empty case\tabularnewline
\midrule
URL & Original address used to validate and acquire the resource. & GitHub, Codeberg, or Zenodo URL\tabularnewline
Normalized identifier & Short platform-specific identifier used for database lookup. & \texttt{owner/notebooks} or \texttt{123456}\tabularnewline
Platform & Source platform detected from the URL or stored database row. & \texttt{github}, \texttt{codeberg}, or \texttt{zenodo}\tabularnewline
Notebook paths & Repository-relative notebook files selected for execution. & \texttt{analysis.ipynb}; empty means discover all notebooks.\tabularnewline
Setup paths & Optional setup files or scripts supplied with the resource. & \texttt{setup.py}; empty means no setup path was supplied.\tabularnewline
Requirement paths & Optional dependency files supplied with the resource. & \texttt{requirements.txt}; empty means infer dependencies from available evidence.\tabularnewline
\bottomrule
\end{tabularx}
\end{table}

These six fields are the handover from input handling to notebook processing. The URL tells the connector what to check and download. The normalized identifier gives the database a short lookup key. The platform prevents confusion between resources with similar names. The three path fields tell the pipeline which notebooks, setup files, or requirement files to use. If a path field is empty, the pipeline discovers the files itself when possible, or continues without that optional input.

After the six fields exist, notebook processing can start. The pipeline then stores its results in two places. The working database is the main evidence: it stores resources, metadata, runs, notebooks, executions, metrics, errors, and classifications. The file system stores supporting artefacts, such as logs, executed notebooks, and JSON comparison files. The knowledge-graph step reads the database, exports selected rows to CSV, and rebuilds the RDF graph from those CSV files and mapping rules. This means the database must be preserved, while the CSV and RDF files can be recreated from it.

\clearpage
\section{Data and Classification Requirements}

The database must separate a resource from the actions performed on it. The same repository or Zenodo record can be processed several times. One run can execute several notebooks, and one notebook can have several classification attempts. If all results were stored in one repository row, a new run would overwrite the old history.

\begin{table}[H]
\centering
\caption{Required relational entities and their purpose.}
\label{tab:data-entities}
\small
\begin{tabularx}{\textwidth}{@{}p{3.6cm}X@{}}
\toprule
Entity & Required meaning\tabularnewline
\midrule
\texttt{repositories} & Stable platform-specific resource identity and supplied path information\tabularnewline
\texttt{repository\_metadata} & One refreshable normalized description linked to the resource\tabularnewline
\texttt{notebooks} & Repository-relative notebook identity and language\tabularnewline
\texttt{repository\_runs} & One processing attempt with status, message, start, end, and duration\tabularnewline
\texttt{notebook\_executions} & Notebook-level execution, error, timing, and cell evidence\tabularnewline
\texttt{notebook\_}\linebreak\texttt{reproducibility\_}\linebreak\texttt{metrics} & Original-versus-executed comparison measurements\tabularnewline
\texttt{notebook\_}\linebreak\texttt{classifications} & Method-specific label, confidence, explanation, model/rule identity, and review state\tabularnewline
\bottomrule
\end{tabularx}
\end{table}

Normalized metadata has one current row for each resource. A later metadata fetch can update that row. Runs, executions, and classifications work differently: every new attempt receives a new row so that earlier results are not deleted. Foreign keys connect the tables, and a uniqueness rule prevents the same notebook path from being added twice to one repository.

The category list must be clear enough for people to use consistently while still covering common notebook purposes. The nine categories are \emph{data preparation}, \emph{data analysis}, \emph{visualization}, \emph{machine learning}, \emph{simulation}, \emph{tutorial}, \emph{software development}, \emph{mixed purpose}, and \emph{uncertain}. Every automated answer must select one of these labels and return a defined data structure, not only free text.

The rules and AI model should examine the same limited information: filename, notebook metadata, selected Markdown and code, imports, cell counts, and output types. Large outputs and possible secrets are not included unless necessary. The pipeline first stores the answer from each method. If the two answers agree, the result can be accepted automatically. Disagreement, an invalid answer, or low confidence is marked for human review. Manual evaluation labels are stored separately so they cannot be confused with predictions.

\clearpage
\section{Evaluation Requirements}

The evaluation must test both the software and the research questions. A demonstration with one repository shows that the route can work, but it is not enough evidence for a thesis comparison. The study also needs connector tests, database checks, classification evaluation, graph tests, and a documented experiment across all three platforms.

\begin{table}[H]
\centering
\caption{Evaluation matrix for the completed system.}
\label{tab:evaluation-matrix}
\small
\begin{tabularx}{\textwidth}{@{}p{3.1cm}p{4.3cm}X@{}}
\toprule
Area & Measures & Required evidence\tabularnewline
\midrule
Connector correctness & validation and acquisition success & Public examples and expected local files\tabularnewline
Metadata normalization & field accuracy and completeness & API-to-database checks for each platform\tabularnewline
Reproducibility route & environment, execution, and output rates & Structured run and notebook records\tabularnewline
Error analysis & counts by stage and category & Logs linked to classified failures\tabularnewline
Classification & accuracy, macro F1, per-class results, agreement & Human-labelled evaluation sample\tabularnewline
Knowledge graph & parseability, expected links, query answers & Triple counts and SPARQL competency tests\tabularnewline
End-to-end behaviour & repeatability and source-DB safety & Clean-run test and file checksums/schema checks\tabularnewline
\bottomrule
\end{tabularx}
\end{table}

The cross-platform experiment must explain how its examples were selected. It must also record the date, software versions, network conditions, and resource identifiers. For each platform, it should report how many resources were selected, validated, and downloaded; which metadata fields were present; how many notebooks were found; and how many environments, executions, and output comparisons succeeded. Percentages must always be shown together with the original counts, especially when the Codeberg and Zenodo samples are smaller than the GitHub data.

Classification must be checked against categories assigned manually without first showing the reviewer the automated answers. Overall accuracy is not enough because some categories may contain far more notebooks than others. The evaluation should also report precision, recall, macro F1, a confusion matrix, and examples of disagreement. The rules, AI model, and combined decision must be measured separately. This shows whether combining them actually improves the result.

Knowledge-graph evaluation must check more than whether the file opens. The ontology and mappings must be valid, CSV row counts must match the database queries, and the graph must contain the expected links between resources, runs, executions, results, and classifications. SPARQL questions test whether these links can be used correctly. A clean rerun must also leave the source database unchanged and recreate all documented outputs. The final claims are therefore supported by stored evidence, not only screenshots.

\chapter{System Design}
\label{chap:design}

\section{Overall Architecture}

The design adds a small connector layer to the original pipeline. GitHub, Codeberg, and Zenodo use different methods for source detection, validation, metadata access, and download. Each connector handles these differences. After the files are stored in a local folder, every platform follows the same notebook-processing route. This avoids repeated code and makes the results easier to compare.

\begin{figure}[ht]
\centering
\fbox{\begin{minipage}{0.92\textwidth}
\centering
\textbf{Input and orchestration}\par
single URL or SQLite batch $\rightarrow$ configuration $\rightarrow$ working database\par
\vspace{0.18cm}
$\downarrow$\par
\vspace{0.12cm}
\textbf{Platform connector boundary}\par
detect $\rightarrow$ validate $\rightarrow$ normalize metadata $\rightarrow$ acquire\par
GitHub $\mid$ Codeberg $\mid$ Zenodo\par
\vspace{0.18cm}
$\downarrow$ local directory and normalized notebook paths\par
\vspace{0.12cm}
\textbf{Shared reproducibility route}\par
discover $\rightarrow$ infer requirements $\rightarrow$ create environment $\rightarrow$ execute $\rightarrow$ compare\par
\vspace{0.18cm}
$\downarrow$\par
\vspace{0.12cm}
\textbf{Evidence and enrichment}\par
SQLite provenance $\rightarrow$ notebook classification $\rightarrow$ CSV/RML $\rightarrow$ RDF graph
\end{minipage}}
\caption{High-level architecture of the cross-platform system.}
\label{fig:overall-architecture}
\end{figure}

The main controller first checks required programs, loads configuration, creates folders, prepares the working database, and selects single or batch mode. It then calls the correct platform connector. The connector produces the common information needed by later steps: repository ID, URL, platform, local folder, notebook paths, and normalized metadata. The remaining functions do not need to know whether the files came from Git or a downloaded archive.

This design also makes future extension easier. A new platform needs a connector that returns the same common information, but it does not need a new environment builder or notebook executor. A new classifier or graph mapping can read stored results without changing the download code. Each part can therefore change without rewriting the whole pipeline.

\clearpage
The layers show where each failure happened. An invalid URL is a connector problem. A failed package installation is an environment problem. An invalid AI response is a classification problem. An invalid mapping is a knowledge-graph problem. Storing the stage gives more useful information than one general ``failed'' status.

\begin{table}[ht]
\centering
\caption{Architectural layers, inputs, and outputs.}
\label{tab:architecture-layers}
\small
\begin{tabularx}{\textwidth}{@{}p{2.6cm}p{3.4cm}p{3.6cm}X@{}}
\toprule
Layer & Main input & Main output & Responsibility\tabularnewline
\midrule
Orchestration & configuration and user choice & prepared run context & dependencies, folders, mode, DB copy\tabularnewline
Connector & URL and platform & local resource and metadata & validate, describe, acquire\tabularnewline
Notebook processing & local files & executions and comparisons & discover, install, run, compare\tabularnewline
Classification & notebook summary & method decisions and review state & assign purpose with provenance\tabularnewline
Persistence & structured events & relational history & stable identity and repeated activities\tabularnewline
Semantic publication & database rows and mappings & RDF graph and query results & interoperable representation\tabularnewline
\bottomrule
\end{tabularx}
\end{table}

The system uses a working database to protect the original data. On the first run, the source database is copied to the output folder. New tables, columns, and results are added only to this copy. Experiments can therefore start again from the same source, and the original data remains separate from newly generated results.

The database keeps the history of repeated work. A repository is stored once, but every processing attempt creates a new run. A notebook is also stored once, while each execution and classification attempt receives its own row. Normalized repository metadata is different: it stores the most recently fetched description together with the fetch time.

Two clear interfaces separate the main parts. First, every connector returns common metadata and a local folder. Second, the knowledge-graph builder reads CSV files with named columns instead of reading shell variables. The connector interface hides platform differences from notebook execution, while the CSV interface separates graph creation from the operational pipeline.

This design supports both a simple demonstration and a larger experiment. One resource can be followed from its URL to its database row, execution, category, and RDF statements. Batch mode uses the same functions for many resources. The demo therefore shows the real system used by the evaluation.

\clearpage
\section{Pipeline Control Flow}

The main controller processes one resource at a time. It first finds or creates the repository ID and then creates a new run row. Even if the next step fails, the database still records which resource was processed and which run failed.

\begin{figure}[ht]
\centering
\fbox{\begin{minipage}{0.91\textwidth}
\centering
\textbf{Resolve resource ID and create RUNNING record}\par
$\downarrow$\par
\textbf{Detect platform} $\rightarrow$ \textbf{platform validator}\par
$\quad$ invalid $\rightarrow$ finalise run and stop\par
$\downarrow$ valid\par
\textbf{Fetch and save eight metadata fields}\par
$\downarrow$\par
\textbf{Clone/pull Git repository or download/extract Zenodo record}\par
$\quad$ acquisition failure $\rightarrow$ finalise run and stop\par
$\downarrow$\par
\textbf{Register/discover notebooks} $\rightarrow$ \textbf{check Python notebooks}\par
$\downarrow$\par
\textbf{requirements} $\rightarrow$ \textbf{environment} $\rightarrow$ \textbf{execute} $\rightarrow$ \textbf{compare}\par
$\downarrow$\par
\textbf{classify, clean up, finalise, and summarise}
\end{minipage}}
\caption{Control flow for one resource. Every terminal branch stores a final status.}
\label{fig:control-flow}
\end{figure}

The source is validated before its metadata is saved. This prevents an unreachable or unsupported URL from receiving a misleading metadata row. GitHub and Codeberg are checked as Git remotes, while Zenodo is checked through its Records API. After validation, the metadata controller calls exactly one platform function and checks that it returned the expected eight fields.

Automatic notebook discovery happens after download because the files must first exist locally. Paths entered by the user can be registered earlier. Before adding a path, the database function checks whether the same notebook already exists, so it does not create duplicates. When no paths are supplied, the pipeline searches the downloaded folder and its subfolders. It then counts all notebooks and Python notebooks. Separate statuses distinguish a resource with no notebooks from one that contains only unsupported notebook languages.

If environment creation, execution, or comparison fails, the current route ends early. Before it ends, the pipeline records the error type and duration, removes the temporary environment when necessary, and saves a final run status. A failed run therefore cannot remain incorrectly marked as \texttt{RUNNING}.

Batch mode reads resources from the database and calls the same controller for each one. It does not use a second, simpler route. If one repository fails, its status is stored and processing continues with the next resource. This is important because failures are part of the research results and should not stop the complete experiment.

\clearpage
\section{Connector Architecture}

Each platform uses a connector with the same general structure. The main controller provides a URL, the platform detector returns \texttt{github}, \texttt{codeberg}, or \texttt{zenodo}, and the selected connector validates the source, fetches metadata, and downloads the content. Although the internal steps differ, every connector returns the same common type of information.

\begin{table}[ht]
\centering
\caption{Connector decisions for the three supported platforms.}
\label{tab:connector-architecture}
\small
\begin{tabularx}{\textwidth}{@{}p{2.2cm}XXX@{}}
\toprule
Decision & GitHub & Codeberg & Zenodo\tabularnewline
\midrule
Detect & host pattern & host pattern & record URL pattern\tabularnewline
Validate & \texttt{git ls-remote} & \texttt{git ls-remote} & HTTP record response and usable files\tabularnewline
Describe & GitHub repository API & Forgejo repository API & Zenodo Records API\tabularnewline
Acquire & clone or pull & clone or pull & select, download, and extract file\tabularnewline
Normalize ID & owner/repository & owner/repository & numeric record ID\tabularnewline
Notebook URL rule & remove \texttt{/blob/branch/} & remove \texttt{/src/branch/} & discover inside archive\tabularnewline
Resource class & Git repository & Git repository & archived research record\tabularnewline
\bottomrule
\end{tabularx}
\end{table}

GitHub and Codeberg both provide a Git remote, so they share clone and update behaviour. They still need different metadata functions because their APIs can organize fields differently. Zenodo uses a separate download process. A Zenodo record can contain several files, so the connector first prefers a source archive and otherwise selects the first supported file. Processing continues only when the record exists and contains a usable download.

The normalized identifier does not include the hostname, but it is always used together with the platform. For example, \texttt{owner/project} on GitHub and the same text on Codeberg are treated as two different resources. The complete original URL is also saved with each run so the exact input remains available.

Connector failures receive clear status values. Examples include an HTTP error, missing Zenodo ID, missing downloadable file, unreachable Git remote, failed clone, or failed archive extraction. These problems are kept separate from notebook execution errors. The evaluation can therefore measure connector success and notebook execution success independently.

The same design can support another platform in the future. A new connector would need platform detection, validation, metadata conversion, download, and a decision about its resource type. If it returns a local folder and the eight common metadata fields, the remaining pipeline can be reused without adding new branches to every later step.

\clearpage
\section{Metadata Model and Normalization Design}

The metadata controller expects exactly eight fields from every platform. Each connector can read a different JSON structure, but it returns the values in the same order. The controller counts them and sends them to one shared database function. The three platforms therefore use one database structure instead of three separate schemas.

\begin{table}[ht]
\centering
\caption{Conceptual mapping into the normalized metadata record.}
\label{tab:metadata-mapping}
\small
\begin{tabularx}{\textwidth}{@{}p{2.3cm}XXX@{}}
\toprule
Common field & GitHub & Codeberg & Zenodo\tabularnewline
\midrule
Title & repository name & repository name & record title\tabularnewline
Description & description & description & record description\tabularnewline
Authors & owner login & owner/user name & joined creator names\tabularnewline
Licence & licence identifier & licence field & record rights/licence ID\tabularnewline
Keywords & topics & topics & record keywords\tabularnewline
DOI & empty unless supplied & empty unless supplied & DOI from record\tabularnewline
Created & \texttt{created\_at} & \texttt{created\_at} & publication/creation date\tabularnewline
Updated & \texttt{updated\_at} & \texttt{updated\_at} & modification date\tabularnewline
\bottomrule
\end{tabularx}
\end{table}

An empty field is allowed when the source provides no value. The connector must not replace a missing DOI with a URL or guess a licence from a filename. This makes the later metadata-completeness results accurate. A future system may add information from another trusted source, but it would need to record where that new information came from.

The Bash functions return one metadata value per line. For this reason, line breaks inside an API description are converted to spaces before the value is returned. The controller stores the lines in an array and rejects an answer that does not contain eight items. Before the SQL statement is created, single quotes in external text are escaped. SQLite itself creates \texttt{fetched\_at}, so an external API cannot choose the local fetch time.

Descriptive information and pipeline activity are stored separately. The metadata table contains what the platform says about the resource. The runs table contains what the pipeline did. Source creation and update dates are therefore not confused with metadata fetch, run start, or run finish times. The platform stays on the repository row because it is part of the resource identity and decides which connector is used.

When metadata is fetched again, an SQLite upsert updates the existing metadata row for that repository instead of creating a duplicate. Earlier processing runs are not overwritten. The database therefore keeps the newest description together with the complete history of reproducibility attempts.

\clearpage
\section{Persistence and Database Design}

SQLite is the shared storage used by the Bash pipeline, classifier, and graph exporter. The database separates resources that are stored once from activities that can happen many times. A second execution therefore creates new evidence instead of overwriting the first result.

\begin{figure}[ht]
\centering
\fbox{\begin{minipage}{0.91\textwidth}
\centering
\texttt{repositories} $1\rightarrow N$ \texttt{notebooks}\qquad
\texttt{repositories} $1\rightarrow N$ \texttt{repository\_runs}\par
\texttt{repositories} $1\rightarrow 1$ \texttt{repository\_metadata}\par
\vspace{0.15cm}
\texttt{repository\_runs} $1\rightarrow N$ \texttt{notebook\_executions}\par
\texttt{notebooks} $1\rightarrow N$ \texttt{notebook\_executions}\par
\vspace{0.15cm}
\texttt{notebook\_executions} $1\rightarrow 0..1$
\texttt{notebook\_reproducibility\_metrics}\par
\texttt{notebooks} $1\rightarrow N$ \texttt{notebook\_classifications}
\end{minipage}}
\caption{Main relational cardinalities. Activity tables preserve repeated attempts.}
\label{fig:database-cardinality}
\end{figure}

The source and working databases have different purposes. The source is copied only when no working database exists. After that, every write uses the working copy. Database updates only add missing tables or columns. Before adding a column, the pipeline checks \texttt{PRAGMA table\_info}. Old repository rows receive GitHub as a compatibility value because the inherited dataset is GitHub-only. The platform of every new resource is detected from its URL.

A repository is identified by its normalized identifier together with its platform. A notebook is unique inside one repository by its relative path. Each repository has one current metadata row. Runs, executions, metrics, and classifications have their own IDs so every attempt can be traced separately. Foreign keys connect these tables.

Statuses use defined values instead of arbitrary summary text. A run starts as \texttt{RUNNING} and later receives a final status, message, and duration. Notebook execution stores the error type, category, message, cell number, and relevant counts. Comparison metrics are written only when both an original and executed notebook are available. Classification rows store the method, category, confidence, explanation, and review information.

This structure supports both detailed inspection and summary statistics. Repository and run rows provide platform success rates. Execution rows provide error counts. The metadata table shows missing fields, and classification links support results by notebook category. The graph exporter also uses the same database IDs to create stable RDF identifiers, so database records and graph resources can be connected clearly.

\clearpage
\section{Notebook Classification Architecture}

Classification is designed as a process whose decisions can be checked. It is not only one unexplained AI request. The process has four stages: read the notebook safely, create a limited summary, run two independent classification methods, and combine their answers with human review when needed.

\begin{figure}[ht]
\centering
\fbox{\begin{minipage}{0.91\textwidth}
\centering
\textbf{Notebook file}\par
$\downarrow$ parse without executing\par
filename $\mid$ metadata $\mid$ Markdown $\mid$ code $\mid$ imports $\mid$ outputs\par
$\swarrow$\hspace{4.8cm}$\searrow$\par
\textbf{Weighted rule classifier}\hspace{1.2cm}\textbf{AI-model classifier}\par
$\searrow$\hspace{4.8cm}$\swarrow$\par
\textbf{Validate labels and compare decisions}\par
$\downarrow$\par
agreement $\rightarrow$ accept\qquad disagreement/low confidence $\rightarrow$ human review
\end{minipage}}
\caption{Hybrid notebook-classification design.}
\label{fig:classification-design}
\end{figure}

The collector reads the notebook JSON without executing any code. It takes selected Markdown and code, records imports and output types, and removes large or unnecessary output content. Limiting the input reduces model cost and the risk of sending secrets or embedded data. The rules and AI model receive the same type of evidence so their results can be compared fairly.

The rule-based method gives points to categories when it finds matching clues. Plotting libraries add points to visualization, machine-learning imports add points to machine learning, and words such as ``tutorial'' or simulation terms support their related categories. The method returns the category with the highest score, a confidence value, and the clues that were matched. Because the rules are saved with a version, a reviewer can see why the category was chosen.

The AI model receives the nine categories, instructions explaining how to choose between them, and the limited notebook summary. It must return one allowed category, a confidence value, and a short explanation in a defined structure. The model provider and model name are stored as configuration. A local or paid model can therefore be tested later without changing the category definitions or database structure.

The pipeline saves both automated answers before choosing the combined result. When the rules and AI model agree, the category can be accepted. When they disagree, return invalid data, or show low confidence, the notebook is marked for review. A human decision is added as a separate record with reviewer information instead of deleting the automated answers. The evaluation can therefore measure the rules, AI model, combined method, and human labels separately.

\clearpage
\section{Knowledge Graph and Provenance Design}

The knowledge graph extends FAIR Jupyter instead of creating a completely new model. The existing ontology already connects publications, software, notebooks, executions, and reproducibility results. New classes and relations add the three platforms, Zenodo records, pipeline runs, metadata, and classifications required by this project.

\begin{figure}[ht]
\centering
\fbox{\begin{minipage}{0.91\textwidth}
\centering
\textbf{Working SQLite database}\par
$\downarrow$ selected queries\par
\textbf{CSV exports with stable IDs and IRI columns}\par
$\downarrow$ YARRRML compiled to RML\par
\textbf{Morph-KGC mapping fragments}\par
$\downarrow$ merge and deduplicate\par
\textbf{Combined N-Triples graph}\par
$\downarrow$ parse checks and SPARQL competency questions
\end{minipage}}
\caption{Regenerable knowledge-graph construction route.}
\label{fig:kg-construction}
\end{figure}

The ontology defines classes for a repository platform, a Git repository resource, and an archived research record. GitHub and Codeberg resources use the Git-repository class, while Zenodo uses the archived-record class. Both can contain notebook resources, but a Zenodo record is not incorrectly described as a Git repository. The graph also represents metadata, repository runs, notebook executions, output-comparison results, and classifications.

PROV-O provides the general structure for processing history. Platforms are represented as agents. Repositories, notebooks, and results are entities. Pipeline runs, notebook executions, graph creation, and classification are activities. Relations state which activity used a resource and which result it generated. These links create a traceable path from an external platform to the final result, while the shared standard makes the graph easier to combine with other systems.

Named SQL queries export selected database data into CSV files. This keeps graph creation separate from changes inside the operational database. Each CSV contains stable IDs and, where useful, complete RDF identifiers. YARRRML provides the readable mapping rules, while the converted RML files are used by Morph-KGC. Each mapping first creates its own N-Triples file, making errors and triple counts easier to locate. The builder then combines the files, removes duplicate triples, and writes a JSON build summary.

SPARQL test questions check whether the graph works as intended. They verify each resource's platform, the separation of Git repositories and Zenodo records, repository-to-notebook links, run-to-execution links, comparison results, and classifications. The graph is accepted only when its files are valid and the queries return the expected patterns. It therefore becomes a usable representation of pipeline evidence, not only a visualization or unchecked export.

\chapter{Cross-Platform Pipeline Implementation}
\label{chap:implementation}

This chapter explains how the design from Chapter~\ref{chap:design} is implemented in the repository. The focus is on executable functions, data passed between them, and the files they produce. Platform concepts that were introduced earlier are not repeated; instead, each section follows the corresponding code path.

\section{Project Structure and Pipeline Entry Point}

The top-level command is \texttt{run.sh}. It uses \texttt{set -e}, checks whether Python, SQLite, \texttt{jq}, \texttt{unzip}, Git, pyenv, Jupyter, and \texttt{nbconvert} are available, and then replaces itself with \texttt{pipeline/main.sh}. Replacing the process with \texttt{exec} keeps the exit status of the main controller visible to the caller. Missing programs are reported before any database or repository is changed.

The main controller resolves \texttt{PROJECT\_ROOT} from its own location and sources the shell modules listed in Table~\ref{tab:implementation-files}. Sourcing is important because the functions share run variables in one shell process. For example, \texttt{process\_repo} exports \texttt{RUN\_ID}, and the Python comparison module reads it from the environment when it inserts execution rows.

\begin{table}[ht]
\centering
\caption{Executable modules and their implementation responsibilities.}
\label{tab:implementation-files}
\small
\begin{tabularx}{\textwidth}{@{}p{3.7cm}X@{}}
\toprule
File & Main responsibility\tabularnewline
\midrule
\texttt{config/config.sh} & Paths, output directories, source and working database, and batch limit\tabularnewline
\texttt{src/db.sh} & Schema creation, migrations, identifiers, and metadata upsert\tabularnewline
\texttt{src/checks.sh} & Git-remote and Zenodo-record validation\tabularnewline
\texttt{src/repo.sh} & Per-resource controller, connectors, metadata adapters, and batch loop\tabularnewline
\texttt{src/requirements.sh} & Requirement-file collection and import discovery\tabularnewline
\texttt{src/pyenv.sh} & Python selection, virtual environment, notebook execution, and cleanup\tabularnewline
\texttt{src/notebooks.sh} & Original/output pairing and comparison dispatch\tabularnewline
\texttt{analysis/} & Notebook diff, execution summary, classification, and persistence\tabularnewline
\bottomrule
\end{tabularx}
\end{table}

After loading these files, \texttt{initialize\_directories} creates the input folder and the output folders for cloned or extracted resources, comparisons, logs, and the working database. \texttt{ensure\_working\_db} copies \texttt{data/db.sqlite} only when \texttt{output/db/db.sqlite} does not yet exist. All later SQL uses \texttt{DB\_FILE}, which points to the working copy. Finally, \texttt{ensure\_pipeline\_tables} creates missing tables and applies compatible column migrations. The controller reaches the mode selection only after this preparation has succeeded.

\clearpage
\section{Run Creation, Single and Batch Processing}

Single mode asks for four values: the resource URL, notebook paths, setup paths, and requirement paths. The last three values use semicolons as separators and may be empty. \texttt{get\_or\_create\_repo\_id} removes the scheme, host, a Zenodo \texttt{record/} or \texttt{records/} segment, and a trailing \texttt{.git}. It returns the existing repository ID or inserts a new row with the detected platform. The main controller exports this ID and calls \texttt{process\_repo} with the four input values.

Repository identity and run identity are deliberately different. The repository row represents the external resource and is reused. \texttt{create\_repository\_run} inserts a new row for every processing attempt:

\begin{lstlisting}[language=bash,caption={Creation of a separate processing attempt.}]
INSERT INTO repository_runs
    (repository_id, url, run_status, started_at)
VALUES
    ($1, '$2', 'RUNNING', datetime('now'));
SELECT last_insert_rowid();
\end{lstlisting}

The returned \texttt{RUN\_ID} connects all notebook executions from this attempt. A repeated run therefore does not overwrite an earlier status or duration. Every terminal path calls \texttt{finalize\_repository\_run}, which replaces \texttt{RUNNING} with a final status, message, finish time, and elapsed seconds.

Batch mode uses the same \texttt{process\_repo} function. \texttt{process\_sqlite\_flow} selects one eligible row at a time, excluding repositories already represented in \texttt{repository\_runs}. It reads the stored identifier and platform, removes CSV quotes and carriage returns, and reconstructs the complete URL. GitHub becomes \texttt{https://github.com/owner/repository}, Codeberg uses the Codeberg host, and Zenodo uses \texttt{https://zenodo.org/records/id}. Unknown platform values are logged and skipped.

The loop retains processed IDs in memory so a problematic row cannot be selected again during the same command. \texttt{TARGET\_COUNT}, with a default of ten, limits the number of completed repository attempts. A failure is stored for the affected row, while the loop continues with the next candidate. At the end, \texttt{print\_run\_summary} queries the database for total, successful, and unsuccessful runs and prints the elapsed time and output locations.

\clearpage
\section{Platform Detection and Validation}

\subsection*{Detection and normalized identity}

\texttt{detect\_platform} performs one host-based decision. A GitHub URL returns \texttt{github}, a Codeberg URL returns \texttt{codeberg}, a Zenodo URL returns \texttt{zenodo}, and every other input returns \texttt{unknown}. The result is calculated once near the beginning of \texttt{process\_repo} and is reused for validation, metadata selection, acquisition, and database lookup.

The original URL is retained on the run row, while the shorter repository identifier is stored on the resource row. Table~\ref{tab:identifier-examples} shows the transformation used by both repository creation and later comparison lookup.

\begin{table}[ht]
\centering
\caption{Examples of platform detection and identifier normalization.}
\label{tab:identifier-examples}
\small
\begin{tabularx}{\textwidth}{@{}p{2.0cm}p{6.8cm}X@{}}
\toprule
Platform & Input & Stored identifier\tabularnewline
\midrule
GitHub & \url{https://github.com/octocat/Hello-World.git} & \texttt{octocat/Hello-World}\tabularnewline
Codeberg & \texttt{https://codeberg.org/user/project} & \texttt{user/project}\tabularnewline
Zenodo & \texttt{https://zenodo.org/records/3362625} & \texttt{3362625}\tabularnewline
\bottomrule
\end{tabularx}
\end{table}

Notebook links are normalized separately because they contain branch-specific segments. For a GitHub link, the code removes the host, owner, repository, \texttt{/blob/}, and branch. For Codeberg it removes the corresponding \texttt{/src/} route. The remaining value is a repository-relative path such as \texttt{notebooks/example.ipynb}. Zenodo paths are not transformed as web links because they are normally discovered after extraction.

Keeping platform and identifier as separate values prevents a host name from being copied into every relation. It also allows batch mode to construct a valid URL without guessing. The complete URL on each run preserves the exact input, including the original host and route.

An unsupported host never falls through to a GitHub default. The detector returns \texttt{unknown}, and batch mode skips that row with its repository ID in the log. In single mode, the same value fails validation instead of being cloned under a misleading platform. This explicit result is important for database safety: a new host cannot silently reuse the API, URL template, or RDF type of one of the supported platforms.

The normalized value is used consistently when a comparison later resolves \texttt{REPO\_ID}. That lookup includes both identifier and platform. Thus, \texttt{user/project} on GitHub and \texttt{user/project} on Codeberg remain distinct even though their host-stripped text is identical. The run URL can still show the exact address used for a particular attempt.

\clearpage
\subsection*{Validation branches and stored outcomes}

GitHub and Codeberg are Git hosts, so \texttt{validate\_repo} uses \texttt{git ls-remote}. This command checks that the URL can be contacted and exposes Git references without downloading the working tree. A successful command returns zero. A failure produces \texttt{INVALID\_REPOSITORY\_URL}, stores the message \texttt{git ls-remote failed}, records the elapsed time, and ends only that repository route.

Zenodo does not provide a Git remote. \texttt{validate\_zenodo} first checks the full URL with a regular expression. The accepted form contains the Zenodo host, \texttt{record} or \texttt{records}, and a numeric identifier. The captured number is then requested from \texttt{/api/records/id} with connection and overall time limits. When \texttt{ZENODO\_TOKEN} is set, the function adds a bearer token without printing it. A malformed, missing, or unreachable record produces \texttt{INVALID\_ZENODO\_RECORD}.

\begin{figure}[ht]
\centering
\fbox{\begin{minipage}{0.90\textwidth}
\centering
\texttt{PLATFORM=github/codeberg} $\rightarrow$ \texttt{git ls-remote}
$\rightarrow$ valid Git resource\par
\vspace{0.18cm}
\texttt{PLATFORM=zenodo} $\rightarrow$ URL pattern $\rightarrow$ Records API
$\rightarrow$ valid archived record\par
\vspace{0.18cm}
any failed check $\rightarrow$ final run status and duration $\rightarrow$ stop this route
\end{minipage}}
\caption{Validation selects a protocol that matches the resource type.}
\label{fig:validation-implementation}
\end{figure}

Validation occurs before metadata fetching. The metadata adapters therefore receive a source that has already passed its platform check. This avoids storing an API error document as repository metadata and gives a precise failure stage. It also avoids unnecessary downloads: a source that cannot be validated never reaches clone or archive extraction.

The validators only test source availability and form. They do not claim that every notebook will execute. Acquisition, notebook discovery, package installation, and execution have their own statuses. This separation lets later SQL queries distinguish a connector problem from a reproducibility problem.

\clearpage
\section{GitHub Baseline Integration}

GitHub provides the baseline route inherited from the original pipeline. The extension keeps the Git operations unchanged and places them behind the same platform decision used by the new connectors. If the local repository folder is missing, \texttt{process\_repo} performs a shallow clone with \texttt{--depth 1}. If the folder already exists, it runs \texttt{git pull}. Both commands write detailed output to the repository log rather than mixing it with the short terminal summary.

The metadata adapter is new to the baseline route. \texttt{fetch\_github\_metadata} requests \texttt{/repos/owner/repository} from the GitHub REST API \cite{githubapi}. It adds the recommended JSON accept header and includes an authorization header only when \texttt{GITHUB\_TOKEN} is set. A missing or empty response returns an adapter failure before any values are saved.

\begin{table}[ht]
\centering
\caption{GitHub API values returned through the common metadata contract.}
\label{tab:github-adapter}
\small
\begin{tabularx}{\textwidth}{@{}p{3.0cm}p{4.3cm}X@{}}
\toprule
Common value & GitHub JSON expression & Treatment\tabularnewline
\midrule
Title & \texttt{name} & empty string when absent\tabularnewline
Description & \texttt{description} & preserved as text\tabularnewline
Authors & \texttt{owner.login} & one repository owner\tabularnewline
Licence & \texttt{license.spdx\_id} & \texttt{NOASSERTION} becomes empty\tabularnewline
Keywords & \texttt{topics} & joined with semicolons\tabularnewline
DOI & no direct field & empty string\tabularnewline
Dates & \texttt{created\_at}, \texttt{updated\_at} & original API timestamps\tabularnewline
\bottomrule
\end{tabularx}
\end{table}

The adapter uses \texttt{jq} to emit exactly one line per value. This form is intentionally identical to the Codeberg and Zenodo adapters, so the controller does not contain GitHub-specific database code. Repository topics are joined into one text field, and a DOI remains empty rather than being guessed from another identifier.

After clone or update, the route is fully shared. GitHub notebooks use the same registration, requirement inference, pyenv environment, \texttt{nbconvert} execution, output comparison, and database tables as the other platforms. This preserved baseline is useful in the evaluation because differences can be attributed to connectors and source data instead of to three separate execution implementations.

\clearpage
\section{Codeberg Connector Implementation}

\subsection*{Reuse of Git behaviour}

Codeberg is implemented as a Git connector rather than as a copy of the GitHub pipeline. Detection selects \texttt{codeberg}, validation still uses \texttt{git ls-remote}, and acquisition still uses shallow clone or pull. The difference appears at boundaries that expose host-specific data: URL reconstruction, full notebook-link normalization, metadata API access, and the platform column.

The metadata function calls the Forgejo repository endpoint at \url{https://codeberg.org/api/v1/repos/owner/repository} \cite{codebergapi,forgejoapi}. It uses an Accept header and optionally sends \texttt{CODEBERG\_TOKEN}. The response is converted into the same eight lines as GitHub. The implementation checks several possible owner properties because Forgejo responses can contain a full name, login, or username.

Licence handling is also defensive. A missing value becomes an empty string. If the value is an object, the adapter tries an SPDX identifier, key, and readable name in that order. If it is another JSON type, it is converted to text. This logic prevents a complete JSON object from being written into the short licence column.

\begin{lstlisting}[language=bash,caption={Selection of a Codeberg owner value in the adapter.}]
if (.owner.full_name // "") != ""
then .owner.full_name
else (.owner.login // .owner.username // "")
end
\end{lstlisting}

Topics are joined with semicolons, and repository creation and update timestamps are retained. Codeberg, like GitHub, has no direct DOI field in this endpoint, so the adapter returns an empty DOI. The absence is stored and can later be measured as metadata completeness.

The connector writes \texttt{codeberg} on the repository row. This value controls URL reconstruction in batch mode and produces a separate platform resource in the knowledge graph. It is not inferred again from a local directory name, because the same owner/repository text could exist on more than one Git host.

\clearpage
\subsection*{Path handling and an end-to-end example}

A Codeberg notebook link normally contains \texttt{/src/branch/} instead of GitHub's \texttt{/blob/branch/}. \texttt{insert\_notebooks\_from\_paths} accepts either full form and applies one regular expression with the alternatives \texttt{blob|src}. Downstream functions consequently receive only the relative file name. This small normalization prevents the environment code from trying to open a web URL below the local repository directory.

Consider the input \texttt{https://codeberg.org/tplasdio/ipynb-py-convert}. The main controller performs the following concrete sequence:

\begin{enumerate}[leftmargin=1.6em]
  \item store \texttt{tplasdio/ipynb-py-convert} with platform \texttt{codeberg};
  \item create a \texttt{RUNNING} row containing the complete Codeberg URL;
  \item validate the remote with \texttt{git ls-remote};
  \item request Forgejo metadata and upsert the eight normalized values;
  \item clone the Git repository or update its existing local folder;
  \item register supplied notebooks or discover all \texttt{.ipynb} files;
  \item execute the shared environment, execution, and comparison route;
  \item store the final status and elapsed time.
\end{enumerate}

No Codeberg branch exists after step five. This is the point at which platform-specific work ends. Requirement paths, setup paths, and notebook paths have the same relative form as GitHub paths and can be processed by the existing functions.

Mixed batch processing uses the stored platform rather than assuming every identifier belongs to GitHub. Values read through SQLite's CSV output are stripped of quotes and Windows carriage returns before the switch reconstructs the URL. If a Codeberg row fails, its run is finalized and the batch advances. A single unavailable forge project therefore does not cancel results for GitHub or Zenodo.

The metadata regression test uses this public Codeberg example together with one GitHub and one Zenodo source. It checks that all three adapters create a metadata row and that every stored title is non-empty. The Codeberg path is therefore exercised through the same controller and persistence contract as the other platforms.

\clearpage
\section{Zenodo Connector and Content Acquisition}

\subsection*{Record validation, metadata, and file selection}

Zenodo is an archived record rather than a Git repository. Its connector therefore replaces the Git operations while preserving the same output contract. \texttt{zenodo\_record\_id} extracts the numeric part of a record URL. Both \texttt{record} and \texttt{records} are accepted, so older and current link forms produce the same identifier.

After validation, \texttt{fetch\_zenodo\_metadata} requests the Records API \cite{zenodoapi}. Title and description come from the record metadata. The creator array is reduced to names and joined with semicolons. Licence data may be a text value or an object; the function selects an identifier or title. Keywords are joined, and the DOI is selected from the direct record field, metadata, or persistent-identifier structure. Creation and modification dates also use ordered fallbacks. Newlines in all eight returned values are replaced by spaces so one metadata field cannot be mistaken for several array elements.

The file list is used separately for acquisition. \texttt{zenodo\_download\_url} reads the JSON response, filters for names ending in \texttt{.zip}, and chooses the first ZIP when one exists. If no ZIP is listed, it selects the first available file. The selected download link is returned to \texttt{fetch\_zenodo}. A record that passes the API check therefore has a clear acquisition decision rather than being sent to Git.

\begin{table}[ht]
\centering
\caption{Zenodo connector values for record \texttt{3362625}.}
\label{tab:zenodo-example}
\small
\begin{tabularx}{\textwidth}{@{}p{3.3cm}X@{}}
\toprule
Item & Connector representation\tabularnewline
\midrule
Input URL & \texttt{https://zenodo.org/records/3362625}\tabularnewline
Stored resource ID & \texttt{3362625}\tabularnewline
Validation & HTTP request to \texttt{/api/records/3362625}\tabularnewline
DOI & \texttt{10.5281/zenodo.3362625}\tabularnewline
Acquisition & API-selected archive download and extraction\tabularnewline
Local result & directory below \texttt{output/cloned\_repos/}\tabularnewline
\bottomrule
\end{tabularx}
\end{table}

The directory name is derived with the same \texttt{basename} step used for Git resources. Because the URL ends in the record number, the Zenodo folder remains stable and can be addressed by all shared modules.

\clearpage
\subsection*{Extraction and convergence with the shared route}

\texttt{fetch\_zenodo} creates the destination, downloads the selected file as a temporary archive, extracts it, and removes the archive after a successful extraction. Some deposits contain one outer directory. The function detects this case and moves its contents one level upward, including hidden files, before removing the empty wrapper. The final directory therefore has the same role as a cloned Git working tree.

This local-directory boundary is what allows Zenodo to use the original notebook-processing code. The next stage searches recursively for \texttt{.ipynb} files when no paths were supplied. Each absolute match is converted to a path relative to the extracted record directory and is passed through the normal notebook insertion function. Requirement and setup files found inside the record are handled exactly like files in a Git checkout.

\begin{figure}[ht]
\centering
\fbox{\begin{minipage}{0.91\textwidth}
\centering
Zenodo URL $\rightarrow$ numeric ID $\rightarrow$ Records API\par
$\downarrow$\par
select archive $\rightarrow$ download $\rightarrow$ extract $\rightarrow$ flatten one wrapper folder\par
$\downarrow$\par
local resource directory $\rightarrow$ discover notebooks $\rightarrow$ shared reproducibility route
\end{minipage}}
\caption{The Zenodo connector converts an archived record into the pipeline's local-directory contract.}
\label{fig:zenodo-convergence}
\end{figure}

The route records separate outcomes for a malformed or missing record, an acquisition that does not create a directory, a record with no notebooks, and a record with no Python notebooks. Each outcome includes its elapsed time. The source can therefore be valid even when its deposited files are unsuitable for the execution stage.

The connector does not describe Zenodo as Git merely because the extracted files resemble a repository. The database platform stays \texttt{zenodo}, the normalized identifier stays numeric, and the graph exporter assigns the archived-record class. These values preserve the nature of the published source after all three platforms have converged on a local folder.

\clearpage
\section{Cross-Platform Metadata Normalization}

\texttt{fetch\_and\_save\_repo\_metadata} is the controller for all descriptive data. It receives the repository ID, original URL, and detected platform. After normalizing the URL to an API path, a \texttt{case} statement calls exactly one adapter. \texttt{mapfile -t} captures the eight output lines in an indexed Bash array.

The controller checks \texttt{\${\#metadata[@]}} before writing. When the count is not eight, it records a clear message and leaves the metadata table unchanged. When the count is correct, indexes zero through seven are passed to \texttt{save\_repo\_metadata} in the declared order. This boundary prevents each connector from creating its own SQL statement.

The database writer applies \texttt{sql\_escape} to every external text value. This helper doubles single quotes, which keeps a title such as \texttt{Researcher's notebook} inside one SQLite string literal. The writer then uses \texttt{ON CONFLICT(repo\_id) DO UPDATE}. The unique repository key produces one current descriptive record, while \texttt{fetched\_at} is refreshed with the local SQLite clock.

\begin{table}[ht]
\centering
\caption{Separation between resource metadata and processing provenance.}
\label{tab:metadata-provenance-storage}
\small
\begin{tabularx}{\textwidth}{@{}p{4.0cm}XX@{}}
\toprule
Stored item & Source & Updated behaviour\tabularnewline
\midrule
Title, description, authors & platform API & refresh current metadata row\tabularnewline
Licence, keywords, DOI & platform API & keep empty when absent\tabularnewline
Created and updated dates & platform API & describe the external resource\tabularnewline
Metadata fetch time & working database clock & refresh on every successful fetch\tabularnewline
Run start, finish, duration & pipeline controller & add a new run row\tabularnewline
\bottomrule
\end{tabularx}
\end{table}

This implementation makes missing data explicit. GitHub and Codeberg normally produce empty DOI values, while Zenodo can provide a DOI. Empty values are not replaced with guesses. Consequently, later completeness queries compare actual source metadata rather than information invented by the pipeline. Similar descriptive elements appear in DataCite's metadata schema \cite{datacite45}, while the dedicated fetch time records the local observation.

\clearpage
\section{Notebook Registration and Discovery}

Notebook paths enter the route in two ways. In single mode, the user can supply relative paths or full GitHub and Codeberg notebook links. \texttt{insert\_notebooks\_from\_paths} splits the semicolon-separated text, trims each item, reduces supported web links to relative paths, and ignores every value that does not end in \texttt{.ipynb}. Before insertion, it queries the repository and path pair. An existing row is reused; a missing row is inserted with the Python language value.

The function stores identity, not notebook content. The \texttt{notebooks} row contains its database ID, repository ID, relative name, and language. The JSON notebook remains in the acquired resource directory and is read only by later Python modules. This keeps the database compact and lets the original and executed files be compared directly.

When \texttt{NOTEBOOK\_PATHS} is empty, discovery waits until clone or extraction has completed. A recursive \texttt{find} collects every \texttt{.ipynb}, sorts the matches, removes the local directory prefix, and joins them with semicolons. The same insertion function then registers the discovered paths. Running the database count before this search would incorrectly classify a new Zenodo record as empty; the implemented order avoids that error.

\begin{figure}[ht]
\centering
\fbox{\begin{minipage}{0.90\textwidth}
\centering
supplied paths $\rightarrow$ normalize $\rightarrow$ register missing rows\par
\textbf{or}\par
empty input $\rightarrow$ acquire files $\rightarrow$ recursive discovery $\rightarrow$ register rows\par
$\downarrow$\par
count all notebooks and Python notebooks $\rightarrow$ continue or store a precise status
\end{minipage}}
\caption{Two notebook-collection routes converge before language checks.}
\label{fig:notebook-collection}
\end{figure}

\texttt{get\_notebook\_language\_stats} returns two values from one SQL query: the total notebook count and the number whose language equals Python without regard to case. A zero total produces \texttt{NO\_NOTEBOOKS}; a nonzero total with zero Python notebooks produces \texttt{NO\_PYTHON\_NOTEBOOKS}. The environment is built only when at least one supported notebook is registered.

\clearpage
\section{Dependency and Environment Reconstruction}

\subsection*{Combining declared and observed requirements}

\texttt{process\_requirements} creates three temporary text files: packages copied from requirement files, imports detected in notebooks, and their combined content. When explicit requirement paths are present, each path is resolved below \texttt{REPO\_DIR}. Existing files are appended, while missing files create warnings and do not stop the route.

For notebook imports, Jupyter converts each selected notebook into a temporary Python file. Lines beginning with \texttt{import} or \texttt{from} are reduced to their top-level module name. The implementation removes duplicates and separates three cases. Standard-library modules such as \texttt{json}, \texttt{pathlib}, and \texttt{sqlite3} are skipped because they should not be installed from PyPI. A module whose \texttt{.py} file exists inside the resource is treated as local and is also skipped. Remaining names are treated as external packages.

\begin{figure}[ht]
\centering
\fbox{\begin{minipage}{0.90\textwidth}
\centering
declared requirement files $\searrow$\par
\hspace{2.7cm} combine $\rightarrow$ remove empty/comment lines $\rightarrow$ sort once\par
notebook imports $\nearrow$\par
$\downarrow$\par
resource-level \texttt{requirements.txt}
\end{minipage}}
\caption{Requirement reconstruction combines explicit files with imports observed in notebooks.}
\label{fig:requirement-reconstruction}
\end{figure}

The final merge removes blank lines and comments, sorts the remaining entries, and writes one \texttt{requirements.txt} in the resource root. Temporary converted Python files and intermediate lists are deleted. If no package is found, an empty requirement file is still created so the environment function receives a predictable path.

Import discovery provides package names but cannot reliably recover versions or operating-system libraries. The implementation therefore prefers versions already declared by the resource and uses imports as additional evidence. Installation results are logged package by package, which makes an unavailable or renamed dependency visible during later error analysis.

\clearpage
\subsection*{Python selection and isolated execution environment}

\texttt{detect\_python\_version} checks version hints in a fixed order: \texttt{binder/runtime.txt}, a root \texttt{runtime.txt}, \texttt{.python-version}, \texttt{python\_requires} in \texttt{setup.py}, and the corresponding field in \texttt{setup.cfg}. If none is present, it selects Python~3.10. The chosen source is written to the log.

\texttt{ensure\_pyenv\_version} first looks for an exact installed version. A major/minor value such as 3.8 is resolved to the newest stable patch shown by pyenv. The function reuses an installed match or installs the resolved interpreter. Failure produces \texttt{PYTHON\_INSTALL\_FAIL}; a missing interpreter binary produces \texttt{PYTHON\_BINARY\_MISSING}. The use of pyenv follows its documented version-management purpose \cite{pyenv}.

\texttt{setup\_pyenv\_env} removes any older environment for the same resource and creates a new \texttt{venv}. It upgrades \texttt{pip}, \texttt{setuptools}, and \texttt{wheel}, then installs Jupyter and \texttt{nbconvert}. Requirement entries are installed one at a time. A failed package is logged and skipped, allowing execution to reveal whether the package was actually needed. Supplied setup paths are installed from their containing directories with \texttt{pip install .}.

\begin{table}[ht]
\centering
\caption{Environment state shared with the notebook executor.}
\label{tab:environment-state}
\small
\begin{tabularx}{\textwidth}{@{}p{3.5cm}X@{}}
\toprule
Variable & Meaning\tabularnewline
\midrule
\texttt{REPO\_VENV\_DIR} & temporary environment dedicated to the current resource\tabularnewline
\texttt{REPO\_PYTHON} & selected interpreter inside that environment\tabularnewline
\texttt{REPO\_PIP} & package installer inside the same environment\tabularnewline
\texttt{ENV\_ERROR\_TYPE} & structured failure category returned to the controller\tabularnewline
\texttt{ENV\_ERROR\_MESSAGE} & readable explanation stored on the run\tabularnewline
\bottomrule
\end{tabularx}
\end{table}

The environment is temporary by design. \texttt{cleanup\_pyenv\_env} removes it after setup failure, execution failure, or completed comparison. A later resource cannot accidentally reuse packages from the previous one, while the logs and database keep the evidence needed to explain the run.

\clearpage
\section{Notebook Execution and Output Comparison}

\texttt{run\_in\_pyenv\_env} splits the normalized path list and executes each existing notebook from its own directory. Running there preserves relative file access used by notebook code. Jupyter \texttt{nbconvert} writes \texttt{name\_output.ipynb} beside the original and uses the Python~3 kernel \cite{nbconvert}. The \texttt{--allow-errors} option preserves a notebook containing a runtime error so the error cell can be analysed instead of losing the complete output document.

For every path, the executor records duration and one of three observable outcomes in \texttt{notebook\_execution\_times.log}: output created without an error, output created with error cells, or execution failure. When no output notebook is created for any path, the repository route receives \texttt{NOTEBOOK\_EXECUTION\_ERROR}. Otherwise, comparison continues for every original/output pair that exists.

\texttt{compare\_notebook\_outputs} resolves repository and notebook IDs, creates a comparison JSON path, and calls \texttt{analysis/compare\_notebook.py}. The Python program loads both notebooks and uses \texttt{nbdime} to obtain a cell-level diff. Execution counts and notebook metadata are ignored because they do not represent result changes. Source and output changes are recorded separately, and output changes are classified as numeric, textual, or structural.

Potential nondeterminism is detected from code patterns such as random-number calls, current time, environment variables, and UUID creation. These signals do not prove that an output changed; they identify cells whose results can vary between runs. Runtime error outputs are read directly from the executed notebook and are grouped into dependency, file, data, code, resource, network, or execution-environment categories.

For a notebook with $n$ code cells and $i$ unchanged cells, the stored score is
\[
  \textit{reproducibility score}=\begin{cases}
  i/n, & n>0,\\
  1, & n=0.
  \end{cases}
\]
The summary inserts one \path{notebook_executions} row and one linked \path{notebook_reproducibility_metrics} row. It also writes the detailed diff to \path{output/comparisons/}. The original notebook is never overwritten.

Execution and comparison use database uniqueness at the run-and-notebook level. A notebook appears once in one repository run, while another run can create a new execution and metrics pair. The JSON file contains the readable summary, raw notebook diff, and extracted runtime errors. The relational rows retain the fields needed for aggregate queries, and the JSON keeps details that would be inconvenient to split across many SQL columns.

An output notebook that was not created receives a small failure JSON document containing its path, repository and notebook IDs, status, and reason. This ensures that the comparison folder also explains missing pairs. Successful and unsuccessful notebook paths can therefore be inspected without relying only on terminal output.

\clearpage
\section{Error Handling, Testing, and Container Verification}

Expected failures are converted into data close to the point where they occur. Connector failures finalize the repository run immediately. Environment failures set \texttt{ENV\_ERROR\_TYPE} and \texttt{ENV\_ERROR\_MESSAGE}. \texttt{analyze\_env\_error} searches the repository log for known patterns, including missing kernels, missing modules, import or syntax errors, memory problems, timeouts, and network failures. The controller stores the selected type and elapsed seconds before cleanup.

\begin{table}[ht]
\centering
\caption{Automated checks implemented for the completed pipeline.}
\label{tab:pipeline-tests}
\small
\begin{tabularx}{\textwidth}{@{}p{4.2cm}X@{}}
\toprule
Test & Behaviour checked\tabularnewline
\midrule
\texttt{test\_metadata.sh} & valid and invalid Zenodo records plus normalized metadata for all three platforms\tabularnewline
\texttt{test\_repo\_discovery.sh} & acquisition precedes automatic notebook discovery and registration\tabularnewline
\texttt{test\_pipeline.sh} & a complete repository attempt creates run evidence\tabularnewline
\texttt{test\_notebook\_}\linebreak\texttt{classification.py} & collection, rules, AI schema, reconciliation, storage, review, and evaluation\tabularnewline
\texttt{test\_kg.sh} & relational export, materialisation, graph structure, and SPARQL queries\tabularnewline
\bottomrule
\end{tabularx}
\end{table}

The metadata test creates a temporary SQLite database, applies the real migrations, and inserts one repository per platform. It validates Zenodo URL handling, calls the three live metadata adapters, and asserts that exactly three metadata rows with non-empty titles are stored. The discovery test avoids network access with local function replacements but runs the real \texttt{process\_repo} order; it proves that an empty path list reaches discovery before notebook counts are checked.

The classification suite uses a generated notebook and a mocked model response so its eight tests are deterministic. It checks AST import collection, capped rule weights, structured AI output, agreement and disagreement, database persistence, human review, and evaluation metrics. The suite completes without requiring an installed model.

The Binder Dockerfile packages the surrounding tools and offers a stable verification environment. A local image build followed by \texttt{tests/test\_metadata.sh} confirms that curl, Git, \texttt{jq}, SQLite, Python, and the project files work together. Per-repository Python environments remain separate from this container: the image verifies the pipeline runtime, while pyenv reconstructs the environment of each analysed notebook.

Together, these checks cover the route from external source to stored result. They also keep failures observable, which is necessary because an unsuccessful notebook is part of the reproducibility evidence rather than merely a test interruption.

\chapter{Notebook Classification}
\label{chap:classification}

The reproducibility route describes whether a notebook can be executed and whether its outputs change. Classification adds a different result: the main purpose of the notebook. The implementation keeps this decision separate from execution, uses two independent methods, and stores enough evidence for later human checking.

\section{Classification Objective and Taxonomy}

The classifier assigns one primary category to each notebook. Categories describe what the notebook mainly does, not the research subject and not whether its execution succeeded. A biomedical visualization and an astronomy visualization therefore share the same purpose category. This distinction allows execution results to be grouped by workflow type across disciplines and platforms.

The category names and definitions are declared once in \texttt{categories.py}. The rule method, AI method, command-line validation, database review, and evaluation all import this shared definition. A label that is not in the taxonomy is rejected before it reaches storage.

\begin{table}[ht]
\centering
\caption{Notebook-purpose taxonomy used by both automated methods.}
\label{tab:classification-taxonomy}
\small
\begin{tabularx}{\textwidth}{@{}p{3.5cm}X@{}}
\toprule
Category & Operational meaning\tabularnewline
\midrule
\texttt{data\_preparation} & cleans, transforms, combines, or prepares data\tabularnewline
\texttt{data\_analysis} & explores data, computes statistics, or answers research questions\tabularnewline
\texttt{visualization} & mainly creates plots, charts, dashboards, or visual reports\tabularnewline
\texttt{machine\_learning} & trains, evaluates, or applies learned predictive models\tabularnewline
\texttt{simulation} & models systems or performs simulated or numerical experiments\tabularnewline
\texttt{tutorial} & teaches a method or demonstrates a tool or dataset\tabularnewline
\texttt{software\_}\linebreak\texttt{development} & develops, tests, or demonstrates reusable software\tabularnewline
\texttt{mixed\_purpose} & has two or more equally central purposes\tabularnewline
\texttt{uncertain} & contains too little evidence for a reliable decision\tabularnewline
\bottomrule
\end{tabularx}
\end{table}

The first seven labels are core categories. \texttt{mixed\_purpose} is used only when strong evidence is divided between purposes, and \texttt{uncertain} is used when evidence is weak. These two labels prevent the system from forcing every notebook into an apparently precise class. The human-review layer can later replace an uncertain or mixed automated result with a final category while keeping the original answer.

\clearpage
\section{Feature Extraction and Input Representation}

\texttt{collect\_notebook} reads the file with \texttt{nbformat}; it never executes notebook code \cite{nbformat}. The collector separates Markdown, code, and raw cells and records the number of each. It also copies notebook metadata, extracts top-level imports, and converts outputs into short summaries. Image bytes, long text outputs, and traceback contents are not included in the classifier input.

Python imports are parsed with the standard-library \texttt{ast} module. For \texttt{import pandas as pd}, the collector keeps \texttt{pandas}; for \texttt{from sklearn.model\_selection import train\_test\_split}, it keeps \texttt{sklearn}. Lines beginning with IPython magics, shell commands, or help syntax are removed before parsing. A cell beginning with a cell magic is skipped because its contents may use another language. If the remaining text is not valid Python, the collector returns no imports for that cell instead of failing the notebook.

Outputs are represented by type. A display or execution result keeps only its MIME types, a stream keeps the stream name, and an error keeps its error name. An output such as a PNG plot therefore contributes \texttt{image/png} without placing the encoded image into the rule engine or AI prompt.

\begin{table}[ht]
\centering
\caption{Collected notebook evidence and its use.}
\label{tab:classification-features}
\small
\begin{tabularx}{\textwidth}{@{}p{3.2cm}p{4.1cm}X@{}}
\toprule
Evidence & Representation & Used for\tabularnewline
\midrule
Metadata & kernel and language dictionaries & context and compatibility\tabularnewline
Markdown & ordered cell text & stated goal, headings, teaching language\tabularnewline
Code & ordered source text & operations, API calls, and model use\tabularnewline
Imports & unique top-level names & package signals\tabularnewline
Outputs & type and MIME summaries & plots, results, and error signals\tabularnewline
Cell counts & Markdown, code, and raw totals & compact structural context\tabularnewline
\bottomrule
\end{tabularx}
\end{table}

The collected object has a version number, currently \texttt{1.0.0}. The final result also contains a SHA-256 hash of the notebook file. The hash distinguishes two versions with the same path and makes it possible to determine which notebook content was classified. If the classifier is tested with an in-memory collected object, the same hash is calculated from its sorted JSON representation.

\clearpage
\section{Rule-Based Notebook Classification}

\subsection{Classification Rules}

The rule classifier is transparent: every point comes from a named signal. Package imports contribute three points to their related category. Examples include \texttt{pandas} for data preparation, \texttt{statsmodels} for data analysis, \texttt{matplotlib} for visualization, \texttt{sklearn} for machine learning, \texttt{simpy} for simulation, and \texttt{pytest} for software development.

Regular expressions search the combined Markdown and code. They cover both descriptive phrases and concrete operations. For example, ``feature engineering'' and \texttt{fillna(} support data preparation; ``hypothesis test'' and \texttt{corr(} support analysis; ``dashboard'' and \texttt{plt.} support visualization; and \texttt{train\_test\_split} or \texttt{.predict(} support machine learning. Tutorial and software-development patterns use similarly direct phrases and calls.

Each matched pattern adds 1.5 points per occurrence, but at most three occurrences are counted. The cap prevents one repeated word from dominating the notebook. The implementation still records the uncapped match count, the number of counted matches, the cells containing the signal, total searched cells, cell coverage, and awarded points. A reviewer can therefore see whether a signal appeared across the notebook or was repeated many times in one cell.

Image outputs provide an additional visualization signal. Up to three outputs containing an \texttt{image/} MIME type add one point each. This feature complements plotting imports: a notebook that stores figures without using a recognized plotting package can still receive visualization evidence.

\begin{lstlisting}[language=Python,caption={Simplified weighted signal update.}]
for pattern in category_patterns:
    matches = find_all(pattern, notebook_text)
    counted = min(len(matches), 3)
    scores[category] += counted * 1.5

if image_outputs:
    scores["visualization"] += min(image_outputs, 3)
\end{lstlisting}

The rule set is versioned as \texttt{1.1.0}. Every stored rule result contains this version, the score per category, the evidence strings, and the detailed pattern statistics. Changing a weight or expression therefore creates a traceable method version instead of silently changing the meaning of earlier results.

\clearpage
\subsection{Rule Priorities and Conflict Resolution}

After all signals are counted, categories are sorted by decreasing score and then alphabetically to make ties deterministic. The two leading values control the result. If the highest score is below two points, the notebook becomes \texttt{uncertain} with confidence zero. This rule prevents a single weak phrase from creating a strong label.

When evidence is sufficient, the leading core category becomes primary. Other categories are retained as secondary when their score is at least two and at least 65\% of the leading score. These secondary labels do not replace the main decision; they show additional substantial purposes.

The confidence value uses both dominance and signal strength. Let $s_1$ be the highest category score and let $S$ be the sum of all category scores. For an ordinary primary category, the implementation calculates
\[
  c=\frac{s_1}{S}\min\left(1,\frac{s_1}{6}\right).
\]
The first part measures how much of the evidence supports the winner. The second part reduces confidence when the absolute evidence is small. The result is rounded to three decimal places.

Mixed purpose has a stricter test. The second score must be at least three points and at least 90\% of the first. Both leading categories then become secondary explanations below the primary label \texttt{mixed\_purpose}. Its confidence uses their combined score and is capped at 0.95. This avoids describing a close competition as a clear single purpose.

\begin{table}[ht]
\centering
\caption{Rule decision cases in order of application.}
\label{tab:rule-decisions}
\small
\begin{tabularx}{\textwidth}{@{}p{3.0cm}p{5.2cm}X@{}}
\toprule
Case & Condition & Result\tabularnewline
\midrule
Insufficient evidence & top score below 2 & \texttt{uncertain}\tabularnewline
Mixed purpose & second score at least 3 and at least 90\% of top & two leading categories retained\tabularnewline
Clear leader & sufficient evidence without mixed condition & top core category\tabularnewline
Secondary purpose & score at least 2 and 65\% of top & recorded as secondary\tabularnewline
\bottomrule
\end{tabularx}
\end{table}

This order makes the decision reproducible. Given the same collected notebook and rule version, the same scores, category, and confidence are produced without network access or random sampling.

\clearpage
\subsection{Rule-Based Output}

The rule method returns a \texttt{ClassifierResult} data object rather than only a label. Its fields include method name, primary category, secondary categories, confidence, a short reason, evidence, all category scores, detailed signal statistics, and rule version. The same result structure is used by the AI method where fields are applicable.

For a notebook containing \texttt{numpy}, \texttt{matplotlib}, a Markdown heading ``Visualization example'', one \texttt{plt.plot} call, and a PNG output, the rules select visualization. The stored evidence can contain \texttt{import:matplotlib}, the visualization text pattern, the plotting-call pattern, and \texttt{image outputs:1}. \texttt{numpy} may also create data-analysis evidence, which remains visible in the complete score dictionary.

\begin{lstlisting}[language=Python,caption={Abbreviated rule output for a plotting notebook.}]
{
  "method": "rule_based",
  "primary_category": "visualization",
  "secondary_categories": ["data_analysis"],
  "confidence": 0.61,
  "evidence": ["import:matplotlib", "image outputs:1"],
  "scores": {
    "visualization": 8.5,
    "data_analysis": 3.0
  },
  "version": "1.1.0"
}
\end{lstlisting}

The exact confidence depends on all matched signals, so the example shows the structure rather than a fixed result for every plotting notebook. Detailed statistics are nested below the category and pattern. They preserve the distinction between a pattern appearing 50 times and receiving points for only three occurrences.

This output supports two uses. First, the reconciliation layer can compare a precise rule category and confidence with the AI answer. Second, a person can audit the label without reading every cell. If an import is misleading, the other scores and cell coverage show how much influence it had.

The method does not claim that high confidence is a calibrated probability. It is a consistent strength measure derived from the implemented rules. Chapter~8 evaluates whether that internal measure corresponds to correct human labels. Until a human review is available, the database clearly identifies the result as rule-based and keeps the full evidence.

\clearpage
\section{AI-Based Notebook Classification}

\subsection{Model Selection}

The second method uses an AI language model through an HTTP interface. The default configuration points to Ollama and the \texttt{gemma3:4b} model. Gemma~3 provides a four-billion-parameter option suitable for local desktop or server execution \cite{gemma3}. The model name and base URL can be replaced through command arguments or environment variables, so the classification contract is not tied to one provider.

The AI method is intentionally independent of the rules. It receives the category definitions and notebook evidence but not the rule scores or chosen category. An agreement therefore represents two methods reaching the same answer rather than the model copying the rule result.

\subsection{Prompt Design}

The prompt begins with the task and the nine allowed definitions. It then gives decision rules: choose mixed purpose only when several goals are equally central, choose uncertain when evidence is insufficient, treat confidence as a self-assessment rather than a calibrated probability, and cite short observable evidence. A security instruction says that notebook text and code are untrusted data and must never be followed as instructions.

The required response is described with JSON Schema. It contains one allowed primary category, a list of allowed secondary categories, a number from zero to one, a reason, and an evidence list. The same schema is passed in the Ollama \texttt{format} field. Ollama documents that this field can constrain a response to a supplied JSON schema \cite{ollamastructured}. Temperature zero is used to reduce unnecessary variation.

\begin{figure}[ht]
\centering
\fbox{\begin{minipage}{0.91\textwidth}
\centering
system role and fixed task\par
$+$ taxonomy and decision rules\par
$+$ JSON output schema\par
$+$ compact notebook evidence treated as untrusted data\par
$\downarrow$\par
structured category, confidence, reason, and evidence
\end{minipage}}
\caption{The AI request separates trusted instructions from untrusted notebook content.}
\label{fig:ai-prompt-structure}
\end{figure}

The model endpoint returns a message containing JSON. The client parses the message, verifies that the primary label belongs to the taxonomy, removes invalid or duplicate secondary labels, clamps confidence to the interval from zero to one, and records the actual model name returned by the server.

\clearpage
\subsection*{Prompt limits and validation}

\subsection{Input Context Construction}

\texttt{build\_llm\_input} turns the collected notebook into a bounded document. It includes the path, cell counts, kernel and language metadata, unique imports, selected Markdown, selected code, and at most 100 output summaries. Markdown is limited to 6,000 characters and code to 8,000 characters. Cells are considered in their original order until the limit is reached.

These limits keep requests predictable and avoid sending embedded binary outputs. They also reduce the influence of very long generated cells. The selected input still contains the notebook's stated purpose, main operations, packages, and output forms. Because the collector does not execute code, classification can run for a notebook whose reproducibility attempt failed.

Notebook text is serialized as JSON below the prompt instructions. A Markdown cell containing ``ignore the categories'' remains data, not a new system command. The instruction to treat content as untrusted is repeated close to the notebook input. This does not remove every model risk, but it provides a clear boundary and prevents the pipeline from intentionally treating notebook prose as control text.

\subsection{Structured AI Output}

The client sends a non-streaming chat request with the schema in both the prompt and \texttt{format}. A valid answer becomes a \texttt{ClassifierResult} with method \texttt{local\_llm}, primary and secondary categories, confidence, reason, evidence, model name, prompt version, and raw response. The prompt version is \texttt{1.0.0}, so results remain interpretable when instructions change.

Network errors, non-success HTTP responses, missing message fields, invalid JSON, unknown labels, and invalid confidence values are converted into \texttt{LocalLLMError}. The raw exception does not end classification. The hybrid controller records that the model was unavailable and sends the notebook to human review.

The structured response solves a practical storage problem. Free prose would require a second parser and could invent category names. Schema validation ensures that every successful AI result fits the same database fields as the rule result. The reason and evidence remain readable, while the primary label and confidence are directly comparable.

\clearpage
\section{Hybrid Classification Strategy}

\texttt{classify\_notebook} collects the notebook once, runs the rules, requests the AI answer, and passes both to \texttt{\_reconcile}. The default confidence threshold is 0.55. The reconciliation function returns an agreement status, a Boolean review flag, a warning, and an optional provisional category.

\begin{table}[ht]
\centering
\caption{Hybrid reconciliation outcomes.}
\label{tab:hybrid-outcomes}
\footnotesize
\begin{tabularx}{\textwidth}{@{}p{3.6cm}p{5.2cm}X@{}}
\toprule
Status & Condition & Provisional decision\tabularnewline
\midrule
\texttt{AGREED} & same label and both confidences meet threshold & shared label, no review\tabularnewline
\texttt{AGREED\_LOW\_CONFIDENCE} & same label but one confidence is low & shared label, review\tabularnewline
\texttt{PARTIAL\_AGREEMENT} & one primary occurs as the other's secondary & none, review\tabularnewline
\texttt{CLASSIFIER\_}\linebreak\texttt{DISAGREEMENT} & confident primary labels differ & none, review\tabularnewline
\texttt{RULE/LLM\_LOW\_}\linebreak\texttt{CONFIDENCE} & one or both methods are weak & none, review\tabularnewline
\texttt{BOTH\_UNCERTAIN} & both primary labels are uncertain & none, review\tabularnewline
\texttt{LLM\_UNAVAILABLE} & request or validation failed & rule label shown, review\tabularnewline
\texttt{LLM\_NOT\_REQUESTED} & rule-only command was selected & rule label shown, review\tabularnewline
\bottomrule
\end{tabularx}
\end{table}

Only confident exact agreement is accepted without review. Agreement is useful but is not treated as proof: the review-queue command can include all rows and draw a blind random sample, including agreed items. This supports evaluation against independent human decisions, following the broader human-in-the-loop principle of keeping people responsible for uncertain automated decisions \cite{mosqueira2023hitl}.

\section{Classification Integration into the Pipeline}

The classification command accepts notebook files or directories. Directories are searched recursively, paths are resolved, and duplicates are removed while order is preserved. Each notebook produces a complete JSON result. With \texttt{--db-file}, the same result is inserted into \texttt{notebook\_classifications}; with \texttt{--output}, all results are also written as a readable JSON document.

Integration occurs through the working database and stable notebook identity. \texttt{resolve\_notebook\_id} compares the stored relative name, POSIX form, and base filename. It links the classification only when exactly one notebook row matches, avoiding a false relation when two repositories contain the same filename.

The command also makes method changes explicit. \texttt{--rule-only} disables the AI request but still creates a review warning. \texttt{--model}, \texttt{--ollama-url}, \texttt{--timeout}, and \texttt{--confidence-threshold} record the experiment configuration at invocation time. A directory command and a single-file command use the same classification function, so their stored result format is identical.

\clearpage
The classification table preserves both automated methods. It stores the notebook hash, rule category and confidence, AI category and confidence, model, prompt version, agreement status, review flag, warning, provisional and final categories, reviewer information, timestamps, and the complete result JSON. A new classification attempt creates a new row, so a later model or prompt does not overwrite an earlier decision.

Human review can be performed for one row or through a CSV queue. The blind queue hides both automated predictions and can select a reproducible random sample using a fixed seed. After reviewers fill the category and their name, \texttt{apply-reviews} validates each label and updates the selected row. The status becomes \texttt{HUMAN\_REVIEWED}, the review flag becomes false, and the automated answers remain unchanged inside their columns and result JSON.

\section{Failure Handling and Classification Examples}

The completed route handles three kinds of classifier failure. Collection failures, such as an unreadable notebook, end that file with a command error. Model failures become \texttt{LLM\_UNAVAILABLE} and retain the rule result for review. Decision uncertainty becomes a normal stored status rather than an exception. These cases can be counted separately during evaluation.

Three short examples illustrate the output:

\begin{itemize}[leftmargin=1.6em]
  \item A plotting notebook imports \texttt{matplotlib}, calls \texttt{plt.plot}, and contains PNG outputs. Rules and AI select visualization above the threshold, producing \texttt{AGREED} and a provisional visualization category.
  \item A notebook is written as a step-by-step lesson but trains a model. The rules select machine learning and the AI selects tutorial with machine learning secondary. The result is \texttt{PARTIAL\_AGREEMENT} and has no provisional category until review.
  \item A short notebook contains only arithmetic and no explanation. Both methods select uncertain, producing \texttt{BOTH\_UNCERTAIN} and a human-review warning.
\end{itemize}

After review, \texttt{evaluate} compares rule, AI, and hybrid provisional labels with final human categories. It reports coverage, accuracy, macro precision, macro recall, macro F1, and per-category support, plus the rule--AI disagreement rate. The implementation has automated tests for collection, weighting, structured model responses, reconciliation, persistence, human review, and these evaluation calculations. Thus, the classification result is not an isolated label: it is a versioned, reviewable activity connected to the notebook and ready for relational and graph-based analysis.

\chapter{Knowledge Graph Extension}
\label{chap:kg}

The operational pipeline stores evidence in SQLite because shell and Python functions can update relational rows directly. The knowledge-graph layer publishes the same evidence as connected RDF statements. It extends the FAIR Jupyter model and can be rebuilt from the working database without manual graph editing.

\section{Relationship to FAIR Jupyter Knowledge Graph}

FAIR Jupyter provides the semantic baseline: the REPRODUCE-ME ontology, mappings for its original publication-centred dataset, and a Morph-KGC materialisation route \cite{samuel2024fairjupyter}. Its existing data covers articles, authors, GitHub repositories, notebooks, code features, requirements, executions, and reproducibility results. Those files remain in \texttt{code/fairjupyter/mapping/} and are not replaced by the extension.

The cross-platform pipeline produces a second, operational data view. Its rows describe GitHub and Codeberg repositories, Zenodo archived records, current metadata, processing attempts, notebook executions, output comparisons, and notebook classifications. Seven new mapping files translate these rows under the same Notebook FAIR namespace.

\begin{table}[ht]
\centering
\caption{Division between the FAIR Jupyter baseline and this extension.}
\label{tab:fairjupyter-relationship}
\small
\begin{tabularx}{\textwidth}{@{}p{3.2cm}XX@{}}
\toprule
Aspect & FAIR Jupyter baseline & Cross-platform extension\tabularnewline
\midrule
Main input & publication/reproducibility CSV data & working pipeline SQLite database\tabularnewline
Source coverage & GitHub-oriented study data & GitHub, Codeberg, and Zenodo\tabularnewline
Resource focus & articles, repositories, notebooks, analyses & platforms, resources, runs, executions, classifications\tabularnewline
Mappings & original YARRRML and RML files & seven separate pipeline mappings\tabularnewline
Output & materialised FAIR Jupyter graph & regenerated pipeline graph using extended ontology\tabularnewline
\bottomrule
\end{tabularx}
\end{table}

The builder does not download a finished online graph and append text to it. It reads the ontology and current CSV data, creates new mapping fragments, and combines them into \texttt{notebookfair-pipeline.nt}. This is a new materialisation that remains compatible with FAIR Jupyter terms.

The original \texttt{reproduce-me.owl} is preserved. New terms are kept in a separate file. This file, \texttt{notebookfair-}\allowbreak\texttt{extension.owl}, imports the baseline ontology. The separation makes the thesis contribution visible and keeps upstream files available for comparison. A later upstream update can be reviewed without mixing it with local platform terms.

\clearpage
\section{Relational Data Export}

\texttt{export\_pipeline\_csv.py} is the bridge between SQLite and the FAIR Jupyter mapping workflow. It resolves the default database as \texttt{output/db/db.sqlite} and writes to \texttt{code/fairjupyter/pipeline\_data/}. The SQLite connection uses a read-only URI, so graph construction cannot change pipeline evidence.

Seven named SQL queries define the export contract. They select columns in a fixed order rather than using \texttt{SELECT *}. The repository query additionally creates a complete platform URL and chooses the RDF resource type. GitHub and Codeberg rows become \texttt{GitRepositoryResource}; Zenodo rows become \texttt{ArchivedResearchRecord}. The classification query calculates \texttt{effective\_category} as the final human category when present, otherwise the provisional category.

\begin{table}[ht]
\centering
\caption{CSV files exported for the pipeline graph.}
\label{tab:kg-csv-exports}
\footnotesize
\begin{tabularx}{\textwidth}{@{}p{5.1cm}p{1.5cm}X@{}}
\toprule
File & Key & Main information\tabularnewline
\midrule
\texttt{repositories.csv} & \texttt{id} & platform, URL, resource type, file counts\tabularnewline
\texttt{repository\_metadata.csv} & \texttt{repo\_id} & title, creators, licence, DOI, dates\tabularnewline
\texttt{notebooks.csv} & \texttt{id} & repository link, path, language\tabularnewline
\texttt{repository\_runs.csv} & \texttt{id} & resource link, status, message, timing\tabularnewline
\texttt{notebook\_executions.csv} & \texttt{id} & run/notebook links, status, cells, errors\tabularnewline
\texttt{notebook\_reproducibility\_}\linebreak\texttt{metrics.csv} & \texttt{id} & execution link, comparison counts, score\tabularnewline
\texttt{notebook\_classifications.csv} & \texttt{id} & method labels, confidence, review, final label\tabularnewline
\bottomrule
\end{tabularx}
\end{table}

\texttt{export\_query} writes the cursor column names as the CSV header and retrieves rows with \texttt{fetchmany(1000)}. Each batch is written immediately and contributes to a printed row count. Memory use therefore depends on the batch size rather than the number of notebooks in the database. If the classification table exists but has no rows, the normal query creates a header-only file. For compatibility with an older database, a missing classification table also creates the expected empty header.

The exporter validates the database path, creates the destination, closes SQLite in a \texttt{finally} block, and returns a nonzero status on file or SQL errors. A successful run prints one count per CSV and a final completion message. These counts later become part of the graph build summary and structural tests.

For example, a repository row with identifier \texttt{octocat/Hello-World} and platform \texttt{github} receives the complete URL \url{https://github.com/octocat/Hello-World} and the Git-resource class IRI. A Zenodo row with identifier \texttt{3362625} receives the record URL and archived-record class. Computing these derived columns once in SQL keeps conditional platform logic out of the mapping files.

CSV remains an intermediate format, not a second authoritative database. Quoted fields preserve commas and line breaks according to the CSV writer, headers define the mapping names, and row order follows the database ID. If the working database changes, the next export replaces the files with a new consistent snapshot.

\clearpage
\section{Ontology Extension}

\subsection*{Classes and resource distinction}

The extension uses \texttt{https://w3id.org/notebookfair\#} as its vocabulary namespace and imports REPRODUCE-ME. Its classes reuse PROV-O for entities, activities, and agents; DOAP for Git repositories; and existing notebook terms where possible \cite{prov2013}. The extension adds only concepts needed by the cross-platform database.

\texttt{RepositoryResource} is the neutral parent for any external source processed by the pipeline and is a subclass of \texttt{prov:Entity}. \texttt{GitRepositoryResource} is both a repository resource and a \texttt{doap:GitRepository}. GitHub and Codeberg use this class. \texttt{ArchivedResearchRecord} is another repository-resource subtype and is declared disjoint with the Git class. Zenodo uses the archived type.

The disjointness is important. After extraction, a Zenodo folder can be processed like a clone, but the published source does not become a Git repository. The RDF type preserves this difference even when both resources contain the same kind of notebook file.

\texttt{RepositoryPlatform} is a \texttt{prov:Agent}. The ontology declares GitHub, Codeberg, and Zenodo as named individuals with stable platform IRIs and lowercase platform names. \texttt{hostedOnPlatform} links each resource to one platform and is defined as a subproperty of \texttt{prov:wasAttributedTo}.

\begin{figure}[ht]
\centering
\fbox{\begin{minipage}{0.91\textwidth}
\centering
\texttt{prov:Entity}\par
$\uparrow$\par
\texttt{RepositoryResource}\par
$\swarrow$\hspace{4.0cm}$\searrow$\par
\texttt{GitRepositoryResource}\hspace{1.0cm}\texttt{ArchivedResearchRecord}\par
GitHub, Codeberg\hspace{3.1cm}Zenodo\par
\vspace{0.22cm}
both link through \texttt{hostedOnPlatform} to a \texttt{RepositoryPlatform}
\end{minipage}}
\caption{Resource classes preserve the difference between Git and archived sources.}
\label{fig:kg-resource-classes}
\end{figure}

\texttt{NotebookResource} specializes the REPRODUCE-ME notebook class. It represents the registered notebook identity rather than an embedded copy of the JSON document. The repository-to-notebook relationship is represented in both directions through \texttt{containsNotebook} and \texttt{isNotebookOf}.

\clearpage
\subsection*{Activities, results, and properties}

\texttt{RepositoryRun} and \texttt{NotebookExecution} are \texttt{prov:Activity} classes. A repository \texttt{hasRun}; the inverse \texttt{runOfRepository} specializes \texttt{prov:used}. A run \texttt{hasExecution}; its inverse \texttt{partOfRun} specializes \texttt{prov:wasInformedBy}. An execution also identifies the notebook it used through \texttt{executedNotebook}.

\texttt{ReproducibilityResult} is a \texttt{prov:Entity}. \texttt{hasReproducibilityResult} links an execution to its metrics, and \texttt{resultOfExecution} specializes \texttt{prov:wasGeneratedBy}. This placement means the score belongs to one execution result, not permanently to a repository.

\texttt{NotebookClassificationActivity} is another PROV activity. It uses a notebook through \texttt{classifiedNotebook}, while \texttt{hasClassification} provides the inverse path. Classification is modelled as an activity because the same notebook can be classified several times with different rules, models, prompts, or human decisions.

\begin{figure}[ht]
\centering
\fbox{\begin{minipage}{0.92\textwidth}
\centering
\texttt{RepositoryPlatform}\par
$\uparrow$ \texttt{hostedOnPlatform}\par
\texttt{RepositoryResource} $\rightarrow$ \texttt{RepositoryRun}
$\rightarrow$ \texttt{NotebookExecution} $\rightarrow$ \texttt{ReproducibilityResult}\par
$\downarrow$ \texttt{containsNotebook}\hspace{3.8cm}$\uparrow$ \texttt{executedNotebook}\par
\texttt{NotebookResource} $\rightarrow$ \texttt{NotebookClassificationActivity}
\end{minipage}}
\caption{Main provenance path represented by the extension.}
\label{fig:kg-provenance-path}
\end{figure}

Datatype properties attach the operational values. Resources have identifiers and counts; notebooks have a language; runs and executions have statuses and durations; executions have cell and error fields; and results have identical, different, nondeterministic, and total cell counts plus the score. Classification activities have rule and AI categories, confidence values, model, prompt version, agreement status, review flag, provisional category, final category, and warning.

Standard terms are used for descriptive metadata. Dublin Core provides title, description, creator, licence, subject, identifier, created, and modified. \texttt{metadataFetchedAt} remains a Notebook FAIR property because it describes the pipeline's local observation. The ontology therefore separates what the external resource says from when the pipeline observed it.

\clearpage
\section{YARRRML and RML Mapping Design}

The readable mapping source is stored as seven \texttt{pipeline\_*.yaml} files. YARRRML expresses each CSV source, subject template, and predicate-object rule in a compact form \cite{heyvaert2018yarrrml}. \texttt{compile\_pipeline\_mappings.py} converts these files into RML Turtle with the pinned \texttt{@rmlio/yarrrml-parser} package \cite{rmlspec}.

Stable subject templates join data created from different CSV files. A repository is always \texttt{https://w3id.org/notebookfair/repository/id}; notebooks, runs, executions, results, and classifications use their table name and database ID. Metadata mapping uses the repository subject directly, so title and DOI enrich the existing repository instead of creating a disconnected metadata node.

\begin{lstlisting}[caption={Abbreviated repository mapping in YARRRML.}]
pipeline_repositories:
  sources:
    - [pipeline_data/repositories.csv~csv]
  s: https://w3id.org/notebookfair/repository/$(id)
  po:
    - [a, $(resource_type)~iri]
    - [schema:url, $(repository_url)~iri]
    - [nf:hostedOnPlatform,
       https://w3id.org/notebookfair/platform/$(platform)~iri]
\end{lstlisting}

The repository mapping also creates platform individuals from the platform column. The notebook mapping creates both \texttt{isNotebookOf} and \texttt{containsNotebook}. Runs, executions, results, and classifications follow the same forward-and-inverse approach. These explicit inverse links make common SPARQL paths short and allow structural tests to count each expected relationship.

Typed numeric values use \texttt{xsd:integer} or \texttt{xsd:decimal}. URLs and resource types are emitted as IRIs, while names, statuses, and messages remain literals. Empty CSV fields do not create meaningful values and are handled by the mapping engine.

The compiler can write current RML files or run with \texttt{--check}. Check mode compiles into a temporary directory and compares every generated file with the committed RML version. A difference means that a readable YARRRML file changed without updating its executable mapping. The build also requires exactly seven pipeline RML files, preventing an unnoticed partial graph.

\clearpage
\section{Knowledge Graph Materialisation}

\subsection*{Preparation and fragment generation}

\texttt{build\_pipeline\_kg.py} controls materialisation. It checks that Morph-KGC can be imported, locates all seven mappings, parses the extension ontology as RDF/XML, and parses every mapping as Turtle. These steps catch syntax problems before a long data conversion starts.

Unless \texttt{--skip-export} is used, the builder calls the CSV exporter with the selected database and pipeline-data directory. It then prepares three output folders below \texttt{output/kg}: \texttt{config}, \texttt{logs}, and \texttt{split\_graph}. Old generated INI, log, and N-Triples files are removed by their exact patterns. Source ontology and mapping files are never deleted.

For each mapping, \texttt{write\_config} creates a Morph-KGC INI file. The configuration specifies one mapping, N-Triples output, one process, logging level, and a dedicated log path. Running mappings separately gives each relational dataset its own fragment and timing. If Morph-KGC fails, the builder reports the mapping name and log rather than continuing with an incomplete graph.

\begin{table}[ht]
\centering
\caption{Generated files during one knowledge-graph build.}
\label{tab:kg-generated-files}
\small
\begin{tabularx}{\textwidth}{@{}p{3.6cm}X@{}}
\toprule
Location & Content\tabularnewline
\midrule
\texttt{pipeline\_data/} & seven current CSV inputs\tabularnewline
\texttt{output/kg/config/} & one Morph-KGC configuration per mapping\tabularnewline
\texttt{output/kg/logs/} & mapping log, standard output, and standard error\tabularnewline
\path{output/kg/split_graph/} & one \texttt{.nt} fragment per mapping\tabularnewline
\texttt{output/kg/} & combined graph and JSON build summary\tabularnewline
\bottomrule
\end{tabularx}
\end{table}

Each successful N-Triples fragment is parsed with RDFLib and counted. Parsing verifies that the mapping engine produced valid RDF, while the count helps locate unexpected growth or an empty source. Mapping duration is measured with a monotonic timer and is stored in the summary.

\clearpage
\subsection*{Combination, summary, and repeatable command}

After all fragments succeed, \texttt{combine\_graphs} creates a new RDFLib graph, loads the extension ontology, and then loads every pipeline fragment. RDFLib treats the graph as a set of triples, so identical statements from different files appear once in the combined output. The result is serialized as \texttt{notebookfair-pipeline.nt}.

The build summary records the UTC build time, source database, ontology, combined graph, total unique triple count, row count for every CSV, and the output, triples, and seconds for every mapping. These values connect a graph artefact to its relational snapshot and make later builds comparable.

The complete route can be run through the wrapper:

\begin{lstlisting}[language=bash,caption={Complete graph build command.}]
bash code/fairjupyter/run_fairjupyter_kg.sh
\end{lstlisting}

The wrapper resolves the project root and chooses Python from the FAIR Jupyter virtual environment on Linux or Windows. It then passes the working database and \texttt{output/kg} directory to the builder. A caller can set \texttt{PYTHON\_BIN} explicitly or pass \texttt{--skip-export} when validating unchanged CSV files.

\begin{figure}[ht]
\centering
\fbox{\begin{minipage}{0.91\textwidth}
\centering
working SQLite DB\par
$\downarrow$ read-only named queries\par
seven CSV files $\rightarrow$ seven RML mappings $\rightarrow$ seven N-Triples fragments\par
$\downarrow$ ontology load and graph-set union\par
combined N-Triples graph $+$ build summary\par
$\downarrow$ RDFLib assertions and SPARQL result CSV files
\end{minipage}}
\caption{Implemented materialisation and validation route.}
\label{fig:kg-materialisation-route}
\end{figure}

Generated files can be removed and recreated from the working database, ontology, and mappings. This is preferable to hand-editing N-Triples because every statement has a documented source column and rule. It also means classification rows added later are published by rerunning the same command, not by changing the graph manually.

\clearpage
\section{Representation of Notebook Classifications}

The seventh mapping publishes each stored classification as a \texttt{Notebook}\allowbreak\texttt{Classification}\allowbreak\texttt{Activity}. Its subject IRI contains the classification row ID, and \texttt{classifiedNotebook} points to the stable notebook IRI. The inverse \texttt{hasClassification} is written from the notebook, allowing queries to start from either side.

The activity retains method-specific results instead of publishing only the effective label. \texttt{ruleCategory} and \texttt{ruleConfidence} preserve the transparent method. \texttt{llmCategory}, \texttt{llmConfidence}, \texttt{classifierModel}, and \texttt{classifierPromptVersion} preserve the AI attempt. \texttt{classificationAgreementStatus}, \texttt{needsHumanReview}, and \texttt{classificationWarning} explain how the two answers were combined.

\begin{table}[ht]
\centering
\caption{Classification columns mapped into RDF properties.}
\label{tab:classification-rdf}
\small
\begin{tabularx}{\textwidth}{@{}p{4.0cm}p{4.0cm}X@{}}
\toprule
Relational column & RDF property & Meaning\tabularnewline
\midrule
\texttt{rule\_category} & \texttt{nf:ruleCategory} & rule decision\tabularnewline
\texttt{llm\_category} & \texttt{nf:llmCategory} & AI decision\tabularnewline
\texttt{llm\_model} & \texttt{nf:classifierModel} & model identity\tabularnewline
\texttt{agreement\_status} & \texttt{nf:classification}\linebreak\texttt{AgreementStatus} & reconciliation outcome\tabularnewline
\texttt{needs\_human\_review} & \texttt{nf:needsHumanReview} & review requirement\tabularnewline
\texttt{provisional\_category} & \texttt{nf:provisional}\linebreak\texttt{Category} & accepted automated label\tabularnewline
\texttt{final\_category} & \texttt{nf:finalCategory} & human-confirmed label\tabularnewline
\bottomrule
\end{tabularx}
\end{table}

The mapping exports only rows connected to a registered \texttt{notebook\_id}. A file that cannot be resolved uniquely can still retain its classification in SQLite and JSON, but it does not create an unsupported RDF relation. Once the notebook identity is resolved, the classification becomes part of the graph on the next build.

This representation supports questions that combine purpose with execution. A query can select visualization notebooks, follow them to repository platform, inspect their execution status, and retrieve their reproducibility score. Because the automated and human values remain separate, the query can also restrict results to human-reviewed classifications or compare the methods.

An abbreviated instance has the following form:

\begin{lstlisting}[caption={Example links for one classification activity.}]
nf-classification:17 a nf:NotebookClassificationActivity ;
  nf:classifiedNotebook nf-notebook:7 ;
  nf:ruleCategory "visualization" ;
  nf:llmCategory "tutorial" ;
  nf:classificationAgreementStatus "CLASSIFIER_DISAGREEMENT" ;
  nf:needsHumanReview true .
\end{lstlisting}

The example keeps disagreement as evidence rather than replacing it with one label. After human review, a final category is added to the same activity representation in the next materialisation. Queries can still retrieve the two original predictions and determine why review was needed.

\clearpage
\section{SPARQL Competency Questions}

Seven SPARQL files express the main questions that the graph must answer. They are executable requirements rather than illustrative text. \texttt{run\_pipeline\_sparql\_tests.py} loads the combined N-Triples graph, runs each query, writes a result CSV, and then checks structural counts against the exported CSV files \cite{sparql2013}.

\begin{table}[ht]
\centering
\caption{Implemented SPARQL competency queries.}
\label{tab:sparql-questions}
\footnotesize
\begin{tabularx}{\textwidth}{@{}p{1.0cm}p{5.5cm}X@{}}
\toprule
No. & Question & Main relations\tabularnewline
\midrule
1 & How many resources are hosted on each platform? & \texttt{hostedOnPlatform}\tabularnewline
2 & Which normalized metadata values belong to each resource? & Dublin Core plus platform\tabularnewline
3 & What path connects repository, notebook, run, execution, and score? & containment and activity links\tabularnewline
4 & How many runs have each final status? & \texttt{runStatus}\tabularnewline
5 & What is the average reproducibility score per platform? & platform, run, execution, result\tabularnewline
6 & Which Zenodo resources are archived records with DOI values? & type, platform, identifier\tabularnewline
7 & What automated, provisional, or final category belongs to a notebook? & classification properties\tabularnewline
\bottomrule
\end{tabularx}
\end{table}

The structural checks are stronger than testing whether the file parses. Entity counts for repositories, notebooks, runs, executions, results, and classifications must equal the corresponding CSV row counts. The graph must contain all three platform names. Every Zenodo resource must be an archived record and must not be a Git repository. A Zenodo DOI must be present.

Relation counts are also checked. Every notebook needs one repository containment link; every run needs a repository link; every execution needs a run link; every metrics row needs an execution-to-result link; and every exported classification needs a notebook link. Required queries must return at least one row. The classification query is allowed to be empty only when its source CSV contains no classification rows.

On success, the test writes \texttt{test-summary.json} with test time, graph path, triple count, entity counts, query row counts, and status \texttt{PASSED}. This file provides machine-readable evidence that the graph structure and competency questions agree with the relational export.

\clearpage
\section{Resulting Knowledge Graph}

The confirmed reproducibility snapshot contains 5,244 repository resources, 27,274 notebooks, 116 repository runs, 444 notebook executions, 444 reproducibility-result rows, and three normalized metadata rows. The materialisation created 222,048 unique triples including ontology statements. The six populated mapping fragments contributed the counts shown in Table~\ref{tab:kg-build-result}; the classification mapping participates in the same build and contributes instance triples whenever classification rows are present.

\begin{table}[ht]
\centering
\caption{Measured fragment sizes in the confirmed graph build.}
\label{tab:kg-build-result}
\small
\begin{tabularx}{\textwidth}{@{}Xr@{}}
\toprule
Fragment & Triples\tabularnewline
\midrule
Notebook executions & 5,976\tabularnewline
Notebook reproducibility metrics & 3,996\tabularnewline
Notebooks and repository links & 163,637\tabularnewline
Repositories and platforms & 47,199\tabularnewline
Repository metadata & 22\tabularnewline
Repository runs & 1,044\tabularnewline
\midrule
Combined unique graph, including ontology & 222,048\tabularnewline
\bottomrule
\end{tabularx}
\end{table}

The difference between fragment totals and the combined graph reflects ontology triples and set-based removal of duplicate statements. The notebook fragment is largest because each notebook produces its type, label, identifier, language, repository link, and inverse containment relation.

The stored SPARQL test summary reports three platform groups, three metadata resources, 50 example execution routes, five run-status groups, one platform group with reproducibility results in that snapshot, and one Zenodo archived record. All entity counts and required links match the CSV source, and the test status is \texttt{PASSED}.

The final graph is useful in two complementary forms. N-Triples provides a simple exchange and loading format, while the query-result CSV files provide compact evidence for inspection and demonstration. More importantly, the graph connects information that was separate in SQLite tables: a platform to its resource, the resource to notebooks and runs, runs to executions, executions to comparison results, and notebooks to their classifications. This connected path is the completed semantic output of the cross-platform pipeline.

The build is repeatable from source evidence. Re-running the exporter and mappings changes instance triples only when the working database changes, while ontology statements and IRI templates stay versioned in the repository. The build and test summaries make that update visible through dates, row counts, fragment counts, and query results. The output is therefore both a graph for analysis and a traceable product of the implemented pipeline.


\backmatter
\bibliographystyle{unsrt}
\bibliography{bibliography}

\end{document}
